%% Generated by Sphinx.
\def\sphinxdocclass{report}
\IfFileExists{luatex85.sty}
 {\RequirePackage{luatex85}}
 {\ifdefined\luatexversion\ifnum\luatexversion>84\relax
  \PackageError{sphinx}
  {** With this LuaTeX (\the\luatexversion),Sphinx requires luatex85.sty **}
  {** Add the LaTeX package luatex85 to your TeX installation, and try again **}
  \endinput\fi\fi}
\documentclass[letterpaper,10pt,english]{sphinxmanual}
\ifdefined\pdfpxdimen
   \let\sphinxpxdimen\pdfpxdimen\else\newdimen\sphinxpxdimen
\fi \sphinxpxdimen=.75bp\relax
\ifdefined\pdfimageresolution
    \pdfimageresolution= \numexpr \dimexpr1in\relax/\sphinxpxdimen\relax
\fi
%% let collapsible pdf bookmarks panel have high depth per default
\PassOptionsToPackage{bookmarksdepth=5}{hyperref}
%% turn off hyperref patch of \index as sphinx.xdy xindy module takes care of
%% suitable \hyperpage mark-up, working around hyperref-xindy incompatibility
\PassOptionsToPackage{hyperindex=false}{hyperref}
%% memoir class requires extra handling
\makeatletter\@ifclassloaded{memoir}
{\ifdefined\memhyperindexfalse\memhyperindexfalse\fi}{}\makeatother

\PassOptionsToPackage{booktabs}{sphinx}
\PassOptionsToPackage{colorrows}{sphinx}

\PassOptionsToPackage{warn}{textcomp}

\catcode`^^^^00a0\active\protected\def^^^^00a0{\leavevmode\nobreak\ }
\usepackage{cmap}
\usepackage{fontspec}
\defaultfontfeatures[\rmfamily,\sffamily,\ttfamily]{}
\usepackage{amsmath,amssymb,amstext}
\usepackage{polyglossia}
\setmainlanguage{english}



\setmainfont{FreeSerif}[
  Extension      = .otf,
  UprightFont    = *,
  ItalicFont     = *Italic,
  BoldFont       = *Bold,
  BoldItalicFont = *BoldItalic
]
\setsansfont{FreeSans}[
  Extension      = .otf,
  UprightFont    = *,
  ItalicFont     = *Oblique,
  BoldFont       = *Bold,
  BoldItalicFont = *BoldOblique,
]
\setmonofont{FreeMono}[Scale=0.9,
  Extension      = .otf,
  UprightFont    = *,
  ItalicFont     = *Oblique,
  BoldFont       = *Bold,
  BoldItalicFont = *BoldOblique,
]



\usepackage[Bjarne]{fncychap}
\usepackage{sphinx}

\fvset{fontsize=auto}
\usepackage{geometry}


% Include hyperref last.
\usepackage{hyperref}
% Fix anchor placement for figures with captions.
\usepackage{hypcap}% it must be loaded after hyperref.
% Set up styles of URL: it should be placed after hyperref.
\urlstyle{same}


\usepackage{sphinxmessages}




\title{QuiverCombinatoricsTools}
\date{Jun 01, 2026}
\release{1.0}
\author{Tudor\sphinxhyphen{}Ioan Caba, Mia Lam, Emanuel Roth}
\newcommand{\sphinxlogo}{\vbox{}}
\renewcommand{\releasename}{Release}
\makeindex
\begin{document}

\pagestyle{empty}
\sphinxmaketitle
\pagestyle{plain}
\sphinxtableofcontents
\pagestyle{normal}
\phantomsection\label{\detokenize{index::doc}}


\sphinxAtStartPar
\sphinxtitleref{QuiverCombinatoricsTools} is a SageMath package that adds combinatorial functions to \sphinxtitleref{QuiverTools} to calculate symplectic leaves of quiver varieties, available here \sphinxhref{https://github.com/QuiverCombinatoricsTools/QuiverCombinatoricsTools}{https://github.com/emanuel\sphinxhyphen{}roth/QuiverCombinatoricsTools}. It adds on to the \sphinxtitleref{QuiverTools} package written by Pieter Belmans, Hans Franzen and Gianni Petrella, as seen here \sphinxurl{https://sage.quiver.tools/} and here \sphinxurl{https://github.com/QuiverTools/QuiverTools}, so consult their documentation when needed.

\sphinxAtStartPar
To install it, make sure you have both \sphinxtitleref{QuiverTools} and \sphinxtitleref{QuiverCombinatoricsTools}

\begin{sphinxVerbatim}[commandchars=\\\{\}]
\PYG{n}{sage} \PYG{o}{\PYGZhy{}}\PYG{o}{\PYGZhy{}}\PYG{n}{pip} \PYG{n}{install} \PYG{n}{git}\PYG{o}{+}\PYG{n}{https}\PYG{p}{:}\PYG{o}{/}\PYG{o}{/}\PYG{n}{github}\PYG{o}{.}\PYG{n}{com}\PYG{o}{/}\PYG{n}{QuiverTools}\PYG{o}{/}\PYG{n}{QuiverTools}\PYG{o}{.}\PYG{n}{git}
\PYG{n}{sage} \PYG{o}{\PYGZhy{}}\PYG{o}{\PYGZhy{}}\PYG{n}{pip} \PYG{n}{install} \PYG{n}{git}\PYG{o}{+}\PYG{n}{https}\PYG{p}{:}\PYG{o}{/}\PYG{o}{/}\PYG{n}{github}\PYG{o}{.}\PYG{n}{com}\PYG{o}{/}\PYG{n}{QuiverCombinatoricsTools}\PYG{o}{/}\PYG{n}{QuiverCombinatoricsTools}\PYG{o}{.}\PYG{n}{git}
\end{sphinxVerbatim}

\sphinxAtStartPar
and then you can simply run

\begin{sphinxVerbatim}[commandchars=\\\{\}]
\PYG{k+kn}{from}\PYG{+w}{ }\PYG{n+nn}{quiver}\PYG{+w}{ }\PYG{k+kn}{import} \PYG{o}{*}
\PYG{k+kn}{from}\PYG{+w}{ }\PYG{n+nn}{quivercombinatorics}\PYG{+w}{ }\PYG{k+kn}{import} \PYG{o}{*}
\end{sphinxVerbatim}

\sphinxAtStartPar
to get started.

\sphinxAtStartPar
\sphinxstylestrong{Authors}
\begin{itemize}
\item {} 
\sphinxAtStartPar
Tudor\sphinxhyphen{}Ioan Caba (University of Edinburgh)

\item {} 
\sphinxAtStartPar
Mia Lam (University of Edinburgh)

\item {} 
\sphinxAtStartPar
Emanuel Roth (University of Edinburgh)

\end{itemize}

\sphinxAtStartPar
We were supervised by Gwyn Bellamy (University of Glasgow), as part of an \sphinxhref{https://www.agq-cdt.org/}{AGQ} computing project.


\chapter{Generating quivers}
\label{\detokenize{index:generating-quivers}}
\sphinxAtStartPar
Here are some functions that help generate quivers to test examples.
\index{quiver\_from\_cartan\_matrix() (in module quivercombinatorics)@\spxentry{quiver\_from\_cartan\_matrix()}\spxextra{in module quivercombinatorics}}

\begin{fulllineitems}
\phantomsection\label{\detokenize{index:quivercombinatorics.quiver_from_cartan_matrix}}
\pysigstartsignatures
\pysiglinewithargsret
{\sphinxcode{\sphinxupquote{quivercombinatorics.}}\sphinxbfcode{\sphinxupquote{quiver\_from\_cartan\_matrix}}}
{\sphinxparam{\DUrole{n}{C}}}
{}
\pysigstopsignatures
\sphinxAtStartPar
Returns the quiver \(Q\) given by a Cartan matrix \(C\)

\sphinxAtStartPar
INPUT:
\begin{itemize}
\item {} 
\sphinxAtStartPar
\sphinxcode{\sphinxupquote{C}} – a square matrix with entries in \(\mathbb{Z}\)

\end{itemize}

\sphinxAtStartPar
OUTPUT: The quiver \(Q\) given by a Cartan matrix \(C\)

\sphinxAtStartPar
EXAMPLE:

\begin{sphinxVerbatim}[commandchars=\\\{\}]
\PYG{n}{sage}\PYG{p}{:} \PYG{k+kn}{from}\PYG{+w}{ }\PYG{n+nn}{quivercombinatorics}\PYG{+w}{ }\PYG{k+kn}{import} \PYG{o}{*}
\PYG{n}{sage}\PYG{p}{:} \PYG{n}{C} \PYG{o}{=} \PYG{p}{[}\PYG{p}{[}\PYG{l+m+mi}{2}\PYG{p}{,} \PYG{o}{\PYGZhy{}}\PYG{l+m+mi}{1}\PYG{p}{,} \PYG{l+m+mi}{0}\PYG{p}{]}\PYG{p}{,} \PYG{p}{[}\PYG{o}{\PYGZhy{}}\PYG{l+m+mi}{1}\PYG{p}{,} \PYG{l+m+mi}{2}\PYG{p}{,} \PYG{o}{\PYGZhy{}}\PYG{l+m+mi}{2}\PYG{p}{]}\PYG{p}{,} \PYG{p}{[}\PYG{l+m+mi}{0}\PYG{p}{,} \PYG{o}{\PYGZhy{}}\PYG{l+m+mi}{2}\PYG{p}{,} \PYG{l+m+mi}{2}\PYG{p}{]}\PYG{p}{]}
\PYG{n}{sage}\PYG{p}{:} \PYG{n}{quiver\PYGZus{}from\PYGZus{}cartan\PYGZus{}matrix}\PYG{p}{(}\PYG{n}{C}\PYG{p}{)}
\PYG{n}{a} \PYG{n}{quiver} \PYG{k}{with} \PYG{l+m+mi}{3} \PYG{n}{vertices} \PYG{o+ow}{and} \PYG{l+m+mi}{3} \PYG{n}{arrows}
\end{sphinxVerbatim}

\end{fulllineitems}

\index{random\_quiver() (in module quivercombinatorics)@\spxentry{random\_quiver()}\spxextra{in module quivercombinatorics}}

\begin{fulllineitems}
\phantomsection\label{\detokenize{index:quivercombinatorics.random_quiver}}
\pysigstartsignatures
\pysiglinewithargsret
{\sphinxcode{\sphinxupquote{quivercombinatorics.}}\sphinxbfcode{\sphinxupquote{random\_quiver}}}
{\sphinxparam{\DUrole{n}{vertices}}\sphinxparamcomma \sphinxparam{\DUrole{n}{max\_arrows\_per\_edge}}}
{}
\pysigstopsignatures
\sphinxAtStartPar
Returns a randomly generated quiver

\sphinxAtStartPar
INPUT:
\begin{itemize}
\item {} 
\sphinxAtStartPar
\sphinxcode{\sphinxupquote{vertices}} – number of vertices you want in the quiver, in \(\mathbb{N}_0\)

\item {} 
\sphinxAtStartPar
\sphinxcode{\sphinxupquote{max\_arrows\_per\_edge}} – maximum of number of arrows you want in the quiver, in \(\mathbb{N}_0\)

\end{itemize}

\sphinxAtStartPar
OUTPUT: A randomly generated quiver with the prescribed vertices and maximum number of edges

\sphinxAtStartPar
EXAMPLES:

\begin{sphinxVerbatim}[commandchars=\\\{\}]
\PYG{n}{sage}\PYG{p}{:} \PYG{k+kn}{from}\PYG{+w}{ }\PYG{n+nn}{quivercombinatorics}\PYG{+w}{ }\PYG{k+kn}{import} \PYG{o}{*}
\PYG{n}{sage}\PYG{p}{:} \PYG{n}{random\PYGZus{}quiver}\PYG{p}{(}\PYG{l+m+mi}{3}\PYG{p}{,} \PYG{l+m+mi}{2}\PYG{p}{)}
\PYG{n}{a} \PYG{n}{quiver} \PYG{k}{with} \PYG{l+m+mi}{3} \PYG{n}{vertices} \PYG{o+ow}{and} \PYG{l+m+mi}{6} \PYG{n}{arrows}

\PYG{n}{sage}\PYG{p}{:} \PYG{k+kn}{from}\PYG{+w}{ }\PYG{n+nn}{quivercombinatorics}\PYG{+w}{ }\PYG{k+kn}{import} \PYG{o}{*}
\PYG{n}{sage}\PYG{p}{:} \PYG{n}{random\PYGZus{}quiver}\PYG{p}{(}\PYG{l+m+mi}{46}\PYG{p}{,} \PYG{l+m+mi}{25}\PYG{p}{)}
\PYG{n}{a} \PYG{n}{quiver} \PYG{k}{with} \PYG{l+m+mi}{46} \PYG{n}{vertices} \PYG{o+ow}{and} \PYG{l+m+mi}{26388} \PYG{n}{arrows}
\end{sphinxVerbatim}

\end{fulllineitems}


\sphinxAtStartPar
We notate quivers by \(Q\), with vertices in \(Q_0\) and edges in \(Q_1\).


\chapter{Constructing \protect\(\Sigma_{\lambda}\protect\)}
\label{\detokenize{index:constructing-sigma-lambda}}
\sphinxAtStartPar
We follow William Crawley\sphinxhyphen{}Boevey in \sphinxhref{https://link.springer.com/article/10.1023/A:1017558904030}{this paper}. Given a quiver \(Q\) with \(n\) vertices, and \({\lambda}\in\mathbb{Z}^n\), \(\Sigma_{\lambda}\) is the set of \(\alpha\in\mathbb{N}^n\) such that \(\alpha\) is a positive root of \(Q\), \(\alpha\cdot\lambda = 0\), and
\begin{equation*}
\begin{split}p(\alpha)>\sum_{t=1}^r p(\beta^{(t)})\end{split}
\end{equation*}
\sphinxAtStartPar
for any decomposition \(\alpha=\beta^{(1)}+\dots+\beta^{(r)}\) with \(r\geq 2\) and \(\beta^{(t)}\) a positive root of \(Q\) with \(\lambda\cdot\beta^{(t)}=0\) for all \sphinxtitleref{t}.
\index{p\_function() (quivercombinatorics.quivercombinatorics.Quiver method)@\spxentry{p\_function()}\spxextra{quivercombinatorics.quivercombinatorics.Quiver method}}

\begin{fulllineitems}
\phantomsection\label{\detokenize{index:quivercombinatorics.quivercombinatorics.Quiver.p_function}}
\pysigstartsignatures
\pysiglinewithargsret
{\sphinxcode{\sphinxupquote{Quiver.}}\sphinxbfcode{\sphinxupquote{p\_function}}}
{\sphinxparam{\DUrole{n}{x}}}
{}
\pysigstopsignatures
\sphinxAtStartPar
Outputs the function \(p(x) = 1 - \frac{1}{2}(x, x)\), where \((x, x)\) is the symmetrized Euler form. Note that \(p(x)\geq 0\) if \(x\) is a root, and \(p(x) = 0\) if and only if \(x\) is a real root

\sphinxAtStartPar
INPUT:
\begin{itemize}
\item {} 
\sphinxAtStartPar
\sphinxcode{\sphinxupquote{x}} – an element of \(\mathbb{Z}Q_0\)

\end{itemize}

\sphinxAtStartPar
OUTPUT: \(1 - \frac{1}{2}(x, x)\)

\sphinxAtStartPar
EXAMPLE:

\begin{sphinxVerbatim}[commandchars=\\\{\}]
\PYG{n}{sage}\PYG{p}{:} \PYG{k+kn}{from}\PYG{+w}{ }\PYG{n+nn}{quivercombinatorics}\PYG{+w}{ }\PYG{k+kn}{import} \PYG{o}{*}
\PYG{n}{sage}\PYG{p}{:} \PYG{n}{Q} \PYG{o}{=} \PYG{n}{Quiver}\PYG{p}{(}\PYG{p}{[}\PYG{p}{[}\PYG{l+m+mi}{0}\PYG{p}{,} \PYG{l+m+mi}{1}\PYG{p}{]}\PYG{p}{,} \PYG{p}{[}\PYG{l+m+mi}{1}\PYG{p}{,} \PYG{l+m+mi}{0}\PYG{p}{]}\PYG{p}{]}\PYG{p}{)}
\PYG{n}{sage}\PYG{p}{:} \PYG{n}{Q}\PYG{o}{.}\PYG{n}{p\PYGZus{}function}\PYG{p}{(}\PYG{p}{(}\PYG{l+m+mi}{2}\PYG{p}{,} \PYG{l+m+mi}{3}\PYG{p}{)}\PYG{p}{)}
\PYG{l+m+mf}{0.0}
\end{sphinxVerbatim}

\end{fulllineitems}


\sphinxAtStartPar
We write \(R_\lambda^+\) to denote the set of positive roots \(\alpha\) with \(\alpha\cdot\lambda=0\) and \(\mathbb{N}R_\lambda^+\) for the set of sums of elements of \(R_\lambda^+\). Using \sphinxhref{https://link.springer.com/article/10.1023/A:1017558904030}{Theorem 5.6 of Crawley\sphinxhyphen{}Boevey’s paper}.

\begin{sphinxadmonition}{note}{Theorem (Crawley\sphinxhyphen{}Boevey, 2001)}

\sphinxAtStartPar
For \(\alpha\in\mathbb{N}^n\), then \(\alpha\in\Sigma_\lambda\) if and only if \(0\neq\alpha\in\mathbb{N}R_\lambda^+\) and \((\beta,\alpha-\beta)\leq -2\) whenever \(\beta\in\mathbb{N}^n\) and \(0<\alpha<\beta\).
\end{sphinxadmonition}
\index{R\_lambda\_plus() (quivercombinatorics.quivercombinatorics.Quiver method)@\spxentry{R\_lambda\_plus()}\spxextra{quivercombinatorics.quivercombinatorics.Quiver method}}

\begin{fulllineitems}
\phantomsection\label{\detokenize{index:quivercombinatorics.quivercombinatorics.Quiver.R_lambda_plus}}
\pysigstartsignatures
\pysiglinewithargsret
{\sphinxcode{\sphinxupquote{Quiver.}}\sphinxbfcode{\sphinxupquote{R\_lambda\_plus}}}
{\sphinxparam{\DUrole{n}{l}}\sphinxparamcomma \sphinxparam{\DUrole{n}{v}}}
{}
\pysigstopsignatures
\sphinxAtStartPar
Returns a list of elements of \(R_\lambda^+\), the set of positive roots \(\alpha\) with \(\alpha\cdot\lambda=0\), up to the upper bound \(v\)

\sphinxAtStartPar
INPUT:
\begin{itemize}
\item {} 
\sphinxAtStartPar
\sphinxcode{\sphinxupquote{l}} – an element of \(\mathbb{Z}Q_0\)

\item {} 
\sphinxAtStartPar
\sphinxcode{\sphinxupquote{v}} – an element of \(\mathbb{N}Q_0\)

\end{itemize}

\sphinxAtStartPar
OUTPUT: A list of elements of \(R_\lambda^+\), where \sphinxcode{\sphinxupquote{l}} is \(\lambda\), the set of positive roots \(\alpha\) with \(\alpha\cdot\lambda=0\), up to the upper bound \(v\)

\sphinxAtStartPar
EXAMPLE:

\begin{sphinxVerbatim}[commandchars=\\\{\}]
\PYG{n}{sage}\PYG{p}{:} \PYG{k+kn}{from}\PYG{+w}{ }\PYG{n+nn}{quivercombinatorics}\PYG{+w}{ }\PYG{k+kn}{import} \PYG{o}{*}
\PYG{n}{sage}\PYG{p}{:} \PYG{n}{Q} \PYG{o}{=} \PYG{n}{Quiver}\PYG{p}{(}\PYG{p}{[}\PYG{p}{[}\PYG{l+m+mi}{0}\PYG{p}{,} \PYG{l+m+mi}{1}\PYG{p}{]}\PYG{p}{,} \PYG{p}{[}\PYG{l+m+mi}{1}\PYG{p}{,} \PYG{l+m+mi}{0}\PYG{p}{]}\PYG{p}{]}\PYG{p}{)}
\PYG{n}{sage}\PYG{p}{:} \PYG{n}{Q}\PYG{o}{.}\PYG{n}{R\PYGZus{}lambda\PYGZus{}plus}\PYG{p}{(}\PYG{p}{(}\PYG{l+m+mi}{1}\PYG{p}{,} \PYG{o}{\PYGZhy{}}\PYG{l+m+mi}{1}\PYG{p}{)}\PYG{p}{,} \PYG{p}{(}\PYG{l+m+mi}{5}\PYG{p}{,} \PYG{l+m+mi}{5}\PYG{p}{)}\PYG{p}{)}
\PYG{p}{[}\PYG{p}{(}\PYG{l+m+mi}{1}\PYG{p}{,} \PYG{l+m+mi}{1}\PYG{p}{)}\PYG{p}{,} \PYG{p}{(}\PYG{l+m+mi}{2}\PYG{p}{,} \PYG{l+m+mi}{2}\PYG{p}{)}\PYG{p}{,} \PYG{p}{(}\PYG{l+m+mi}{3}\PYG{p}{,} \PYG{l+m+mi}{3}\PYG{p}{)}\PYG{p}{,} \PYG{p}{(}\PYG{l+m+mi}{4}\PYG{p}{,} \PYG{l+m+mi}{4}\PYG{p}{)}\PYG{p}{]}
\end{sphinxVerbatim}

\end{fulllineitems}


\sphinxAtStartPar
We also need the helper function \sphinxcode{\sphinxupquote{N\_set}}.
\index{N\_set() (in module quivercombinatorics)@\spxentry{N\_set()}\spxextra{in module quivercombinatorics}}

\begin{fulllineitems}
\phantomsection\label{\detokenize{index:quivercombinatorics.N_set}}
\pysigstartsignatures
\pysiglinewithargsret
{\sphinxcode{\sphinxupquote{quivercombinatorics.}}\sphinxbfcode{\sphinxupquote{N\_set}}}
{\sphinxparam{\DUrole{n}{S}}\sphinxparamcomma \sphinxparam{\DUrole{n}{v}}}
{}
\pysigstopsignatures
\sphinxAtStartPar
For a list of vectors \(S\) in \(\mathbb{Z}Q_0\), returns all possible sums of vectors in \(S\) as a list, up to the upper bound \(v\). Though not directly about quivers, this is a helper function to define \(\Sigma_\lambda\)

\sphinxAtStartPar
INPUT:
\begin{itemize}
\item {} 
\sphinxAtStartPar
\sphinxcode{\sphinxupquote{S}} – a list of vectors \(S\) in \(\mathbb{Z}Q_0\)

\item {} 
\sphinxAtStartPar
\sphinxcode{\sphinxupquote{v}} – vector in \(\mathbb{Z}Q_0\)

\end{itemize}

\sphinxAtStartPar
OUTPUT: A list of all possible sums of vectors in \(S\) as a list, up to the upper bound \(v\)

\sphinxAtStartPar
EXAMPLE:

\begin{sphinxVerbatim}[commandchars=\\\{\}]
\PYG{n}{sage}\PYG{p}{:} \PYG{k+kn}{from}\PYG{+w}{ }\PYG{n+nn}{quivercombinatorics}\PYG{+w}{ }\PYG{k+kn}{import} \PYG{o}{*}
\PYG{n}{sage}\PYG{p}{:} \PYG{n}{N\PYGZus{}set}\PYG{p}{(}\PYG{p}{[}\PYG{p}{[}\PYG{l+m+mi}{0}\PYG{p}{,} \PYG{l+m+mi}{1}\PYG{p}{]}\PYG{p}{]}\PYG{p}{,} \PYG{p}{[}\PYG{l+m+mi}{0}\PYG{p}{,} \PYG{l+m+mi}{3}\PYG{p}{]}\PYG{p}{)}
\PYG{p}{[}\PYG{p}{(}\PYG{l+m+mi}{0}\PYG{p}{,} \PYG{l+m+mi}{1}\PYG{p}{)}\PYG{p}{,} \PYG{p}{(}\PYG{l+m+mi}{0}\PYG{p}{,} \PYG{l+m+mi}{2}\PYG{p}{)}\PYG{p}{,} \PYG{p}{(}\PYG{l+m+mi}{0}\PYG{p}{,} \PYG{l+m+mi}{3}\PYG{p}{)}\PYG{p}{]}
\end{sphinxVerbatim}

\end{fulllineitems}


\sphinxAtStartPar
With these functions, we can define \(\Sigma_{\lambda}\).
\index{sigma\_lambda() (quivercombinatorics.quivercombinatorics.Quiver method)@\spxentry{sigma\_lambda()}\spxextra{quivercombinatorics.quivercombinatorics.Quiver method}}

\begin{fulllineitems}
\phantomsection\label{\detokenize{index:quivercombinatorics.quivercombinatorics.Quiver.sigma_lambda}}
\pysigstartsignatures
\pysiglinewithargsret
{\sphinxcode{\sphinxupquote{Quiver.}}\sphinxbfcode{\sphinxupquote{sigma\_lambda}}}
{\sphinxparam{\DUrole{n}{l}}\sphinxparamcomma \sphinxparam{\DUrole{n}{v}}}
{}
\pysigstopsignatures
\sphinxAtStartPar
Returns a list of elements of \(\Sigma_\lambda\), up to the upper bound \(v\)

\sphinxAtStartPar
INPUT:
\begin{itemize}
\item {} 
\sphinxAtStartPar
\sphinxcode{\sphinxupquote{l}} – an element of \(\mathbb{Z}Q_0\)

\item {} 
\sphinxAtStartPar
\sphinxcode{\sphinxupquote{v}} – an element of \(\mathbb{N}Q_0\)

\end{itemize}

\sphinxAtStartPar
OUTPUT: A list of elements of \(\Sigma_\lambda\), where \sphinxcode{\sphinxupquote{l}} is \(\lambda\), up to the upper bound \(v\)

\sphinxAtStartPar
EXAMPLE:

\begin{sphinxVerbatim}[commandchars=\\\{\}]
\PYG{n}{sage}\PYG{p}{:} \PYG{k+kn}{from}\PYG{+w}{ }\PYG{n+nn}{quivercombinatorics}\PYG{+w}{ }\PYG{k+kn}{import} \PYG{o}{*}
\PYG{n}{sage}\PYG{p}{:} \PYG{n}{Q} \PYG{o}{=} \PYG{n}{Quiver}\PYG{p}{(}\PYG{p}{[}\PYG{p}{[}\PYG{l+m+mi}{0}\PYG{p}{,} \PYG{l+m+mi}{1}\PYG{p}{]}\PYG{p}{,} \PYG{p}{[}\PYG{l+m+mi}{1}\PYG{p}{,} \PYG{l+m+mi}{0}\PYG{p}{]}\PYG{p}{]}\PYG{p}{)}
\PYG{n}{sage}\PYG{p}{:} \PYG{n}{Q}\PYG{o}{.}\PYG{n}{sigma\PYGZus{}lambda}\PYG{p}{(}\PYG{p}{(}\PYG{l+m+mi}{1}\PYG{p}{,} \PYG{o}{\PYGZhy{}}\PYG{l+m+mi}{1}\PYG{p}{)}\PYG{p}{,} \PYG{p}{(}\PYG{l+m+mi}{5}\PYG{p}{,} \PYG{l+m+mi}{5}\PYG{p}{)}\PYG{p}{)}
\PYG{p}{[}\PYG{p}{(}\PYG{l+m+mi}{1}\PYG{p}{,} \PYG{l+m+mi}{1}\PYG{p}{)}\PYG{p}{]}
\end{sphinxVerbatim}

\end{fulllineitems}



\chapter{CB\sphinxhyphen{}decompositions}
\label{\detokenize{index:cb-decompositions}}
\sphinxAtStartPar
A \sphinxstyleemphasis{representation type} of \(x\) of a quiver \(Q\) is a sequence \(\tau=(\beta^{(1)},n_1;\dots;\beta^{(k)},n_k)\) with \(\beta^{(i)}\in\Sigma_\lambda\) and \(n_i\in\mathbb{N}\) such that
\begin{equation*}
\begin{split}x=\sum_{i=1}^k n_i\beta^{(i)}\end{split}
\end{equation*}
\sphinxAtStartPar
We allow \(\beta^{(i)}\) to occur multiple times if it is an imaginary root, i.e. \(p(\beta^{(i)})>0\). The \sphinxstyleemphasis{CB\sphinxhyphen{}decomposition} of \(x\) is defined to be the unique representation type of \(x\) for which
\begin{equation*}
\begin{split}p(\tau):=\sum_{i=1}^k p(\beta^{(i)})\end{split}
\end{equation*}
\sphinxAtStartPar
is maximal. In order to construct CB\sphinxhyphen{}decompositions (called canonical decompositions in \sphinxhref{https://link.springer.com/article/10.1023/A:1017558904030}{this paper}, but are \sphinxstyleemphasis{not} the same as \sphinxcode{\sphinxupquote{canonical\_decomposition}} from \sphinxtitleref{QuiverTools}). We first need to find all representation types of \(x\), up to a bound \(v\), for which we need the following helper functions.
\index{vector\_decomposition() (in module quivercombinatorics)@\spxentry{vector\_decomposition()}\spxextra{in module quivercombinatorics}}

\begin{fulllineitems}
\phantomsection\label{\detokenize{index:quivercombinatorics.vector_decomposition}}
\pysigstartsignatures
\pysiglinewithargsret
{\sphinxcode{\sphinxupquote{quivercombinatorics.}}\sphinxbfcode{\sphinxupquote{vector\_decomposition}}}
{\sphinxparam{\DUrole{n}{x}}\sphinxparamcomma \sphinxparam{\DUrole{n}{S}}}
{}
\pysigstopsignatures
\sphinxAtStartPar
For a list of vectors \(S\) in \(\mathbb{Z}Q_0\), returns all possible sums of vectors in \(S\) as a list that sum up to \(x\). This is a helper function in order to define CB\sphinxhyphen{}decompositions

\sphinxAtStartPar
INPUT:
\begin{itemize}
\item {} 
\sphinxAtStartPar
\sphinxcode{\sphinxupquote{x}} – a vector in \(\mathbb{Z}Q_0\)

\item {} 
\sphinxAtStartPar
\sphinxcode{\sphinxupquote{S}} – a list of vectors \(S\) in \(\mathbb{Z}Q_0\)

\end{itemize}

\sphinxAtStartPar
OUTPUT: All possible sums of vectors in \(S\) as a list that sum up to \(x\)

\sphinxAtStartPar
EXAMPLE:

\begin{sphinxVerbatim}[commandchars=\\\{\}]
\PYG{n}{sage}\PYG{p}{:} \PYG{k+kn}{from}\PYG{+w}{ }\PYG{n+nn}{quivercombinatorics}\PYG{+w}{ }\PYG{k+kn}{import} \PYG{o}{*}
\PYG{n}{sage}\PYG{p}{:} \PYG{n}{vector\PYGZus{}decomposition}\PYG{p}{(}\PYG{p}{(}\PYG{l+m+mi}{4}\PYG{p}{,} \PYG{l+m+mi}{5}\PYG{p}{)}\PYG{p}{,} \PYG{p}{[}\PYG{p}{(}\PYG{l+m+mi}{0}\PYG{p}{,} \PYG{l+m+mi}{1}\PYG{p}{)}\PYG{p}{,} \PYG{p}{(}\PYG{l+m+mi}{1}\PYG{p}{,} \PYG{l+m+mi}{0}\PYG{p}{)}\PYG{p}{,} \PYG{p}{(}\PYG{l+m+mi}{1}\PYG{p}{,} \PYG{l+m+mi}{1}\PYG{p}{)}\PYG{p}{]}\PYG{p}{)}
\PYG{p}{[}\PYG{p}{[}\PYG{p}{[}\PYG{p}{(}\PYG{l+m+mi}{0}\PYG{p}{,} \PYG{l+m+mi}{1}\PYG{p}{)}\PYG{p}{,} \PYG{l+m+mi}{1}\PYG{p}{]}\PYG{p}{,} \PYG{p}{[}\PYG{p}{(}\PYG{l+m+mi}{1}\PYG{p}{,} \PYG{l+m+mi}{1}\PYG{p}{)}\PYG{p}{,} \PYG{l+m+mi}{4}\PYG{p}{]}\PYG{p}{]}\PYG{p}{,}
\PYG{p}{[}\PYG{p}{[}\PYG{p}{(}\PYG{l+m+mi}{0}\PYG{p}{,} \PYG{l+m+mi}{1}\PYG{p}{)}\PYG{p}{,} \PYG{l+m+mi}{2}\PYG{p}{]}\PYG{p}{,} \PYG{p}{[}\PYG{p}{(}\PYG{l+m+mi}{1}\PYG{p}{,} \PYG{l+m+mi}{0}\PYG{p}{)}\PYG{p}{,} \PYG{l+m+mi}{1}\PYG{p}{]}\PYG{p}{,} \PYG{p}{[}\PYG{p}{(}\PYG{l+m+mi}{1}\PYG{p}{,} \PYG{l+m+mi}{1}\PYG{p}{)}\PYG{p}{,} \PYG{l+m+mi}{3}\PYG{p}{]}\PYG{p}{]}\PYG{p}{,}
\PYG{p}{[}\PYG{p}{[}\PYG{p}{(}\PYG{l+m+mi}{0}\PYG{p}{,} \PYG{l+m+mi}{1}\PYG{p}{)}\PYG{p}{,} \PYG{l+m+mi}{3}\PYG{p}{]}\PYG{p}{,} \PYG{p}{[}\PYG{p}{(}\PYG{l+m+mi}{1}\PYG{p}{,} \PYG{l+m+mi}{0}\PYG{p}{)}\PYG{p}{,} \PYG{l+m+mi}{2}\PYG{p}{]}\PYG{p}{,} \PYG{p}{[}\PYG{p}{(}\PYG{l+m+mi}{1}\PYG{p}{,} \PYG{l+m+mi}{1}\PYG{p}{)}\PYG{p}{,} \PYG{l+m+mi}{2}\PYG{p}{]}\PYG{p}{]}\PYG{p}{,}
\PYG{p}{[}\PYG{p}{[}\PYG{p}{(}\PYG{l+m+mi}{0}\PYG{p}{,} \PYG{l+m+mi}{1}\PYG{p}{)}\PYG{p}{,} \PYG{l+m+mi}{4}\PYG{p}{]}\PYG{p}{,} \PYG{p}{[}\PYG{p}{(}\PYG{l+m+mi}{1}\PYG{p}{,} \PYG{l+m+mi}{0}\PYG{p}{)}\PYG{p}{,} \PYG{l+m+mi}{3}\PYG{p}{]}\PYG{p}{,} \PYG{p}{[}\PYG{p}{(}\PYG{l+m+mi}{1}\PYG{p}{,} \PYG{l+m+mi}{1}\PYG{p}{)}\PYG{p}{,} \PYG{l+m+mi}{1}\PYG{p}{]}\PYG{p}{]}\PYG{p}{,}
\PYG{p}{[}\PYG{p}{[}\PYG{p}{(}\PYG{l+m+mi}{0}\PYG{p}{,} \PYG{l+m+mi}{1}\PYG{p}{)}\PYG{p}{,} \PYG{l+m+mi}{5}\PYG{p}{]}\PYG{p}{,} \PYG{p}{[}\PYG{p}{(}\PYG{l+m+mi}{1}\PYG{p}{,} \PYG{l+m+mi}{0}\PYG{p}{)}\PYG{p}{,} \PYG{l+m+mi}{4}\PYG{p}{]}\PYG{p}{]}\PYG{p}{]}
\end{sphinxVerbatim}

\end{fulllineitems}

\index{small\_decomposition() (in module quivercombinatorics)@\spxentry{small\_decomposition()}\spxextra{in module quivercombinatorics}}

\begin{fulllineitems}
\phantomsection\label{\detokenize{index:quivercombinatorics.small_decomposition}}
\pysigstartsignatures
\pysiglinewithargsret
{\sphinxcode{\sphinxupquote{quivercombinatorics.}}\sphinxbfcode{\sphinxupquote{small\_decomposition}}}
{\sphinxparam{\DUrole{n}{v}}\sphinxparamcomma \sphinxparam{\DUrole{n}{n}}}
{}
\pysigstopsignatures
\sphinxAtStartPar
For a vector \(v\) in \(\mathbb{Z}Q_0\), it returns all possible partitions of the representation type \([v,n]\)

\sphinxAtStartPar
INPUT:
\begin{itemize}
\item {} 
\sphinxAtStartPar
\sphinxcode{\sphinxupquote{v}} – a vector in \(\mathbb{Z}Q_0\)

\item {} 
\sphinxAtStartPar
\sphinxcode{\sphinxupquote{n}} – a natural number in \(\mathbb{N}\)

\end{itemize}

\sphinxAtStartPar
OUTPUT: All possible partitions of the representation type \([v,n]\)

\sphinxAtStartPar
EXAMPLE:

\begin{sphinxVerbatim}[commandchars=\\\{\}]
\PYG{n}{sage}\PYG{p}{:} \PYG{k+kn}{from}\PYG{+w}{ }\PYG{n+nn}{quivercombinatorics}\PYG{+w}{ }\PYG{k+kn}{import} \PYG{o}{*}
\PYG{n}{sage}\PYG{p}{:} \PYG{n}{small\PYGZus{}decomposition}\PYG{p}{(}\PYG{p}{(}\PYG{l+m+mi}{1}\PYG{p}{,} \PYG{l+m+mi}{1}\PYG{p}{)}\PYG{p}{,} \PYG{l+m+mi}{3}\PYG{p}{)}
\PYG{p}{[}\PYG{p}{[}\PYG{p}{[}\PYG{p}{(}\PYG{l+m+mi}{1}\PYG{p}{,} \PYG{l+m+mi}{1}\PYG{p}{)}\PYG{p}{,} \PYG{l+m+mi}{1}\PYG{p}{]}\PYG{p}{,} \PYG{p}{[}\PYG{p}{(}\PYG{l+m+mi}{1}\PYG{p}{,} \PYG{l+m+mi}{1}\PYG{p}{)}\PYG{p}{,} \PYG{l+m+mi}{1}\PYG{p}{]}\PYG{p}{,} \PYG{p}{[}\PYG{p}{(}\PYG{l+m+mi}{1}\PYG{p}{,} \PYG{l+m+mi}{1}\PYG{p}{)}\PYG{p}{,} \PYG{l+m+mi}{1}\PYG{p}{]}\PYG{p}{]}\PYG{p}{,}
\PYG{p}{[}\PYG{p}{[}\PYG{p}{(}\PYG{l+m+mi}{1}\PYG{p}{,} \PYG{l+m+mi}{1}\PYG{p}{)}\PYG{p}{,} \PYG{l+m+mi}{2}\PYG{p}{]}\PYG{p}{,} \PYG{p}{[}\PYG{p}{(}\PYG{l+m+mi}{1}\PYG{p}{,} \PYG{l+m+mi}{1}\PYG{p}{)}\PYG{p}{,} \PYG{l+m+mi}{1}\PYG{p}{]}\PYG{p}{]}\PYG{p}{,}
\PYG{p}{[}\PYG{p}{[}\PYG{p}{(}\PYG{l+m+mi}{1}\PYG{p}{,} \PYG{l+m+mi}{1}\PYG{p}{)}\PYG{p}{,} \PYG{l+m+mi}{3}\PYG{p}{]}\PYG{p}{]}\PYG{p}{]}
\end{sphinxVerbatim}

\end{fulllineitems}


\sphinxAtStartPar
With these functions, we can determine all representation types.
\index{all\_representation\_types() (quivercombinatorics.quivercombinatorics.Quiver method)@\spxentry{all\_representation\_types()}\spxextra{quivercombinatorics.quivercombinatorics.Quiver method}}

\begin{fulllineitems}
\phantomsection\label{\detokenize{index:quivercombinatorics.quivercombinatorics.Quiver.all_representation_types}}
\pysigstartsignatures
\pysiglinewithargsret
{\sphinxcode{\sphinxupquote{Quiver.}}\sphinxbfcode{\sphinxupquote{all\_representation\_types}}}
{\sphinxparam{\DUrole{n}{l}}\sphinxparamcomma \sphinxparam{\DUrole{n}{x}}}
{}
\pysigstopsignatures
\sphinxAtStartPar
Returns a list of representation types of a quiver, with respect to \(x\)

\sphinxAtStartPar
INPUT:
\begin{itemize}
\item {} 
\sphinxAtStartPar
\sphinxcode{\sphinxupquote{l}} – an element of \(\mathbb{Z}Q_0\)

\item {} 
\sphinxAtStartPar
\sphinxcode{\sphinxupquote{x}} – an element of \(\mathbb{N}Q_0\)

\end{itemize}

\sphinxAtStartPar
OUTPUT: A list of representation types of a quiver, with respect to \(x\), and where \sphinxcode{\sphinxupquote{l}} is \(\lambda\). Each representation type is stored as a list whose elements are 2\sphinxhyphen{}tuples, the first element is in \(\mathbb{N}Q_0\), and the second is in \(\mathbb{N}\)

\sphinxAtStartPar
EXAMPLES:

\begin{sphinxVerbatim}[commandchars=\\\{\}]
\PYG{n}{sage}\PYG{p}{:} \PYG{k+kn}{from}\PYG{+w}{ }\PYG{n+nn}{quivercombinatorics}\PYG{+w}{ }\PYG{k+kn}{import} \PYG{o}{*}
\PYG{n}{sage}\PYG{p}{:} \PYG{n}{Q} \PYG{o}{=} \PYG{n}{Quiver}\PYG{p}{(}\PYG{p}{[}\PYG{p}{[}\PYG{l+m+mi}{0}\PYG{p}{,} \PYG{l+m+mi}{1}\PYG{p}{]}\PYG{p}{,} \PYG{p}{[}\PYG{l+m+mi}{1}\PYG{p}{,} \PYG{l+m+mi}{0}\PYG{p}{]}\PYG{p}{]}\PYG{p}{)}
\PYG{n}{sage}\PYG{p}{:} \PYG{n}{Q}\PYG{o}{.}\PYG{n}{all\PYGZus{}representation\PYGZus{}types}\PYG{p}{(}\PYG{p}{(}\PYG{l+m+mi}{1}\PYG{p}{,} \PYG{o}{\PYGZhy{}}\PYG{l+m+mi}{1}\PYG{p}{)}\PYG{p}{,} \PYG{p}{(}\PYG{l+m+mi}{5}\PYG{p}{,} \PYG{l+m+mi}{5}\PYG{p}{)}\PYG{p}{)}
\PYG{p}{[}\PYG{p}{[}\PYG{p}{[}\PYG{p}{(}\PYG{l+m+mi}{1}\PYG{p}{,} \PYG{l+m+mi}{1}\PYG{p}{)}\PYG{p}{,} \PYG{l+m+mi}{1}\PYG{p}{]}\PYG{p}{,} \PYG{p}{[}\PYG{p}{(}\PYG{l+m+mi}{1}\PYG{p}{,} \PYG{l+m+mi}{1}\PYG{p}{)}\PYG{p}{,} \PYG{l+m+mi}{1}\PYG{p}{]}\PYG{p}{,} \PYG{p}{[}\PYG{p}{(}\PYG{l+m+mi}{1}\PYG{p}{,} \PYG{l+m+mi}{1}\PYG{p}{)}\PYG{p}{,} \PYG{l+m+mi}{1}\PYG{p}{]}\PYG{p}{,} \PYG{p}{[}\PYG{p}{(}\PYG{l+m+mi}{1}\PYG{p}{,} \PYG{l+m+mi}{1}\PYG{p}{)}\PYG{p}{,} \PYG{l+m+mi}{1}\PYG{p}{]}\PYG{p}{,} \PYG{p}{[}\PYG{p}{(}\PYG{l+m+mi}{1}\PYG{p}{,} \PYG{l+m+mi}{1}\PYG{p}{)}\PYG{p}{,} \PYG{l+m+mi}{1}\PYG{p}{]}\PYG{p}{]}\PYG{p}{,}
\PYG{p}{[}\PYG{p}{[}\PYG{p}{(}\PYG{l+m+mi}{1}\PYG{p}{,} \PYG{l+m+mi}{1}\PYG{p}{)}\PYG{p}{,} \PYG{l+m+mi}{1}\PYG{p}{]}\PYG{p}{,} \PYG{p}{[}\PYG{p}{(}\PYG{l+m+mi}{1}\PYG{p}{,} \PYG{l+m+mi}{1}\PYG{p}{)}\PYG{p}{,} \PYG{l+m+mi}{1}\PYG{p}{]}\PYG{p}{,} \PYG{p}{[}\PYG{p}{(}\PYG{l+m+mi}{1}\PYG{p}{,} \PYG{l+m+mi}{1}\PYG{p}{)}\PYG{p}{,} \PYG{l+m+mi}{1}\PYG{p}{]}\PYG{p}{,} \PYG{p}{[}\PYG{p}{(}\PYG{l+m+mi}{1}\PYG{p}{,} \PYG{l+m+mi}{1}\PYG{p}{)}\PYG{p}{,} \PYG{l+m+mi}{2}\PYG{p}{]}\PYG{p}{]}\PYG{p}{,}
\PYG{p}{[}\PYG{p}{[}\PYG{p}{(}\PYG{l+m+mi}{1}\PYG{p}{,} \PYG{l+m+mi}{1}\PYG{p}{)}\PYG{p}{,} \PYG{l+m+mi}{1}\PYG{p}{]}\PYG{p}{,} \PYG{p}{[}\PYG{p}{(}\PYG{l+m+mi}{1}\PYG{p}{,} \PYG{l+m+mi}{1}\PYG{p}{)}\PYG{p}{,} \PYG{l+m+mi}{1}\PYG{p}{]}\PYG{p}{,} \PYG{p}{[}\PYG{p}{(}\PYG{l+m+mi}{1}\PYG{p}{,} \PYG{l+m+mi}{1}\PYG{p}{)}\PYG{p}{,} \PYG{l+m+mi}{3}\PYG{p}{]}\PYG{p}{]}\PYG{p}{,}
\PYG{p}{[}\PYG{p}{[}\PYG{p}{(}\PYG{l+m+mi}{1}\PYG{p}{,} \PYG{l+m+mi}{1}\PYG{p}{)}\PYG{p}{,} \PYG{l+m+mi}{1}\PYG{p}{]}\PYG{p}{,} \PYG{p}{[}\PYG{p}{(}\PYG{l+m+mi}{1}\PYG{p}{,} \PYG{l+m+mi}{1}\PYG{p}{)}\PYG{p}{,} \PYG{l+m+mi}{2}\PYG{p}{]}\PYG{p}{,} \PYG{p}{[}\PYG{p}{(}\PYG{l+m+mi}{1}\PYG{p}{,} \PYG{l+m+mi}{1}\PYG{p}{)}\PYG{p}{,} \PYG{l+m+mi}{2}\PYG{p}{]}\PYG{p}{]}\PYG{p}{,}
\PYG{p}{[}\PYG{p}{[}\PYG{p}{(}\PYG{l+m+mi}{1}\PYG{p}{,} \PYG{l+m+mi}{1}\PYG{p}{)}\PYG{p}{,} \PYG{l+m+mi}{1}\PYG{p}{]}\PYG{p}{,} \PYG{p}{[}\PYG{p}{(}\PYG{l+m+mi}{1}\PYG{p}{,} \PYG{l+m+mi}{1}\PYG{p}{)}\PYG{p}{,} \PYG{l+m+mi}{4}\PYG{p}{]}\PYG{p}{]}\PYG{p}{,}
\PYG{p}{[}\PYG{p}{[}\PYG{p}{(}\PYG{l+m+mi}{1}\PYG{p}{,} \PYG{l+m+mi}{1}\PYG{p}{)}\PYG{p}{,} \PYG{l+m+mi}{2}\PYG{p}{]}\PYG{p}{,} \PYG{p}{[}\PYG{p}{(}\PYG{l+m+mi}{1}\PYG{p}{,} \PYG{l+m+mi}{1}\PYG{p}{)}\PYG{p}{,} \PYG{l+m+mi}{3}\PYG{p}{]}\PYG{p}{]}\PYG{p}{,}
\PYG{p}{[}\PYG{p}{[}\PYG{p}{(}\PYG{l+m+mi}{1}\PYG{p}{,} \PYG{l+m+mi}{1}\PYG{p}{)}\PYG{p}{,} \PYG{l+m+mi}{5}\PYG{p}{]}\PYG{p}{]}\PYG{p}{]}

\PYG{n}{sage}\PYG{p}{:} \PYG{k+kn}{from}\PYG{+w}{ }\PYG{n+nn}{quivercombinatorics}\PYG{+w}{ }\PYG{k+kn}{import} \PYG{o}{*}
\PYG{n}{sage}\PYG{p}{:} \PYG{n}{C} \PYG{o}{=} \PYG{p}{[}\PYG{p}{[}\PYG{l+m+mi}{2}\PYG{p}{,} \PYG{o}{\PYGZhy{}}\PYG{l+m+mi}{1}\PYG{p}{,} \PYG{l+m+mi}{0}\PYG{p}{]}\PYG{p}{,} \PYG{p}{[}\PYG{o}{\PYGZhy{}}\PYG{l+m+mi}{1}\PYG{p}{,} \PYG{l+m+mi}{2}\PYG{p}{,} \PYG{o}{\PYGZhy{}}\PYG{l+m+mi}{2}\PYG{p}{]}\PYG{p}{,} \PYG{p}{[}\PYG{l+m+mi}{0}\PYG{p}{,} \PYG{o}{\PYGZhy{}}\PYG{l+m+mi}{2}\PYG{p}{,} \PYG{l+m+mi}{2}\PYG{p}{]}\PYG{p}{]}
\PYG{n}{sage}\PYG{p}{:} \PYG{n}{Q} \PYG{o}{=} \PYG{n}{quiver\PYGZus{}from\PYGZus{}cartan\PYGZus{}matrix}\PYG{p}{(}\PYG{n}{C}\PYG{p}{)}
\PYG{n}{sage}\PYG{p}{:} \PYG{n}{Q}\PYG{o}{.}\PYG{n}{all\PYGZus{}representation\PYGZus{}types}\PYG{p}{(}\PYG{p}{(}\PYG{l+m+mi}{0}\PYG{p}{,} \PYG{l+m+mi}{0}\PYG{p}{,} \PYG{l+m+mi}{0}\PYG{p}{)}\PYG{p}{,} \PYG{p}{(}\PYG{l+m+mi}{2}\PYG{p}{,} \PYG{l+m+mi}{4}\PYG{p}{,} \PYG{l+m+mi}{3}\PYG{p}{)}\PYG{p}{)}
\PYG{p}{[}\PYG{p}{[}\PYG{p}{[}\PYG{p}{(}\PYG{l+m+mi}{0}\PYG{p}{,} \PYG{l+m+mi}{0}\PYG{p}{,} \PYG{l+m+mi}{1}\PYG{p}{)}\PYG{p}{,} \PYG{l+m+mi}{1}\PYG{p}{]}\PYG{p}{,}
\PYG{p}{[}\PYG{p}{(}\PYG{l+m+mi}{0}\PYG{p}{,} \PYG{l+m+mi}{1}\PYG{p}{,} \PYG{l+m+mi}{0}\PYG{p}{)}\PYG{p}{,} \PYG{l+m+mi}{2}\PYG{p}{]}\PYG{p}{,}
\PYG{p}{[}\PYG{p}{(}\PYG{l+m+mi}{0}\PYG{p}{,} \PYG{l+m+mi}{1}\PYG{p}{,} \PYG{l+m+mi}{1}\PYG{p}{)}\PYG{p}{,} \PYG{l+m+mi}{1}\PYG{p}{]}\PYG{p}{,}
\PYG{p}{[}\PYG{p}{(}\PYG{l+m+mi}{0}\PYG{p}{,} \PYG{l+m+mi}{1}\PYG{p}{,} \PYG{l+m+mi}{1}\PYG{p}{)}\PYG{p}{,} \PYG{l+m+mi}{1}\PYG{p}{]}\PYG{p}{,}
\PYG{p}{[}\PYG{p}{(}\PYG{l+m+mi}{1}\PYG{p}{,} \PYG{l+m+mi}{0}\PYG{p}{,} \PYG{l+m+mi}{0}\PYG{p}{)}\PYG{p}{,} \PYG{l+m+mi}{2}\PYG{p}{]}\PYG{p}{]}\PYG{p}{,}
\PYG{p}{[}\PYG{p}{[}\PYG{p}{(}\PYG{l+m+mi}{0}\PYG{p}{,} \PYG{l+m+mi}{0}\PYG{p}{,} \PYG{l+m+mi}{1}\PYG{p}{)}\PYG{p}{,} \PYG{l+m+mi}{1}\PYG{p}{]}\PYG{p}{,} \PYG{p}{[}\PYG{p}{(}\PYG{l+m+mi}{0}\PYG{p}{,} \PYG{l+m+mi}{1}\PYG{p}{,} \PYG{l+m+mi}{0}\PYG{p}{)}\PYG{p}{,} \PYG{l+m+mi}{2}\PYG{p}{]}\PYG{p}{,} \PYG{p}{[}\PYG{p}{(}\PYG{l+m+mi}{0}\PYG{p}{,} \PYG{l+m+mi}{1}\PYG{p}{,} \PYG{l+m+mi}{1}\PYG{p}{)}\PYG{p}{,} \PYG{l+m+mi}{2}\PYG{p}{]}\PYG{p}{,} \PYG{p}{[}\PYG{p}{(}\PYG{l+m+mi}{1}\PYG{p}{,} \PYG{l+m+mi}{0}\PYG{p}{,} \PYG{l+m+mi}{0}\PYG{p}{)}\PYG{p}{,} \PYG{l+m+mi}{2}\PYG{p}{]}\PYG{p}{]}\PYG{p}{,}
\PYG{p}{[}\PYG{p}{[}\PYG{p}{(}\PYG{l+m+mi}{0}\PYG{p}{,} \PYG{l+m+mi}{0}\PYG{p}{,} \PYG{l+m+mi}{1}\PYG{p}{)}\PYG{p}{,} \PYG{l+m+mi}{2}\PYG{p}{]}\PYG{p}{,} \PYG{p}{[}\PYG{p}{(}\PYG{l+m+mi}{0}\PYG{p}{,} \PYG{l+m+mi}{1}\PYG{p}{,} \PYG{l+m+mi}{0}\PYG{p}{)}\PYG{p}{,} \PYG{l+m+mi}{3}\PYG{p}{]}\PYG{p}{,} \PYG{p}{[}\PYG{p}{(}\PYG{l+m+mi}{0}\PYG{p}{,} \PYG{l+m+mi}{1}\PYG{p}{,} \PYG{l+m+mi}{1}\PYG{p}{)}\PYG{p}{,} \PYG{l+m+mi}{1}\PYG{p}{]}\PYG{p}{,} \PYG{p}{[}\PYG{p}{(}\PYG{l+m+mi}{1}\PYG{p}{,} \PYG{l+m+mi}{0}\PYG{p}{,} \PYG{l+m+mi}{0}\PYG{p}{)}\PYG{p}{,} \PYG{l+m+mi}{2}\PYG{p}{]}\PYG{p}{]}\PYG{p}{,}
\PYG{p}{[}\PYG{p}{[}\PYG{p}{(}\PYG{l+m+mi}{0}\PYG{p}{,} \PYG{l+m+mi}{0}\PYG{p}{,} \PYG{l+m+mi}{1}\PYG{p}{)}\PYG{p}{,} \PYG{l+m+mi}{3}\PYG{p}{]}\PYG{p}{,} \PYG{p}{[}\PYG{p}{(}\PYG{l+m+mi}{0}\PYG{p}{,} \PYG{l+m+mi}{1}\PYG{p}{,} \PYG{l+m+mi}{0}\PYG{p}{)}\PYG{p}{,} \PYG{l+m+mi}{4}\PYG{p}{]}\PYG{p}{,} \PYG{p}{[}\PYG{p}{(}\PYG{l+m+mi}{1}\PYG{p}{,} \PYG{l+m+mi}{0}\PYG{p}{,} \PYG{l+m+mi}{0}\PYG{p}{)}\PYG{p}{,} \PYG{l+m+mi}{2}\PYG{p}{]}\PYG{p}{]}\PYG{p}{,}
\PYG{p}{[}\PYG{p}{[}\PYG{p}{(}\PYG{l+m+mi}{0}\PYG{p}{,} \PYG{l+m+mi}{1}\PYG{p}{,} \PYG{l+m+mi}{0}\PYG{p}{)}\PYG{p}{,} \PYG{l+m+mi}{1}\PYG{p}{]}\PYG{p}{,}
\PYG{p}{[}\PYG{p}{(}\PYG{l+m+mi}{0}\PYG{p}{,} \PYG{l+m+mi}{1}\PYG{p}{,} \PYG{l+m+mi}{1}\PYG{p}{)}\PYG{p}{,} \PYG{l+m+mi}{1}\PYG{p}{]}\PYG{p}{,}
\PYG{p}{[}\PYG{p}{(}\PYG{l+m+mi}{0}\PYG{p}{,} \PYG{l+m+mi}{1}\PYG{p}{,} \PYG{l+m+mi}{1}\PYG{p}{)}\PYG{p}{,} \PYG{l+m+mi}{1}\PYG{p}{]}\PYG{p}{,}
\PYG{p}{[}\PYG{p}{(}\PYG{l+m+mi}{0}\PYG{p}{,} \PYG{l+m+mi}{1}\PYG{p}{,} \PYG{l+m+mi}{1}\PYG{p}{)}\PYG{p}{,} \PYG{l+m+mi}{1}\PYG{p}{]}\PYG{p}{,}
\PYG{p}{[}\PYG{p}{(}\PYG{l+m+mi}{1}\PYG{p}{,} \PYG{l+m+mi}{0}\PYG{p}{,} \PYG{l+m+mi}{0}\PYG{p}{)}\PYG{p}{,} \PYG{l+m+mi}{2}\PYG{p}{]}\PYG{p}{]}\PYG{p}{,}
\PYG{p}{[}\PYG{p}{[}\PYG{p}{(}\PYG{l+m+mi}{0}\PYG{p}{,} \PYG{l+m+mi}{1}\PYG{p}{,} \PYG{l+m+mi}{0}\PYG{p}{)}\PYG{p}{,} \PYG{l+m+mi}{1}\PYG{p}{]}\PYG{p}{,} \PYG{p}{[}\PYG{p}{(}\PYG{l+m+mi}{0}\PYG{p}{,} \PYG{l+m+mi}{1}\PYG{p}{,} \PYG{l+m+mi}{1}\PYG{p}{)}\PYG{p}{,} \PYG{l+m+mi}{1}\PYG{p}{]}\PYG{p}{,} \PYG{p}{[}\PYG{p}{(}\PYG{l+m+mi}{0}\PYG{p}{,} \PYG{l+m+mi}{1}\PYG{p}{,} \PYG{l+m+mi}{1}\PYG{p}{)}\PYG{p}{,} \PYG{l+m+mi}{2}\PYG{p}{]}\PYG{p}{,} \PYG{p}{[}\PYG{p}{(}\PYG{l+m+mi}{1}\PYG{p}{,} \PYG{l+m+mi}{0}\PYG{p}{,} \PYG{l+m+mi}{0}\PYG{p}{)}\PYG{p}{,} \PYG{l+m+mi}{2}\PYG{p}{]}\PYG{p}{]}\PYG{p}{,}
\PYG{p}{[}\PYG{p}{[}\PYG{p}{(}\PYG{l+m+mi}{0}\PYG{p}{,} \PYG{l+m+mi}{1}\PYG{p}{,} \PYG{l+m+mi}{0}\PYG{p}{)}\PYG{p}{,} \PYG{l+m+mi}{1}\PYG{p}{]}\PYG{p}{,} \PYG{p}{[}\PYG{p}{(}\PYG{l+m+mi}{0}\PYG{p}{,} \PYG{l+m+mi}{1}\PYG{p}{,} \PYG{l+m+mi}{1}\PYG{p}{)}\PYG{p}{,} \PYG{l+m+mi}{3}\PYG{p}{]}\PYG{p}{,} \PYG{p}{[}\PYG{p}{(}\PYG{l+m+mi}{1}\PYG{p}{,} \PYG{l+m+mi}{0}\PYG{p}{,} \PYG{l+m+mi}{0}\PYG{p}{)}\PYG{p}{,} \PYG{l+m+mi}{2}\PYG{p}{]}\PYG{p}{]}\PYG{p}{,}
\PYG{p}{[}\PYG{p}{[}\PYG{p}{(}\PYG{l+m+mi}{2}\PYG{p}{,} \PYG{l+m+mi}{4}\PYG{p}{,} \PYG{l+m+mi}{3}\PYG{p}{)}\PYG{p}{,} \PYG{l+m+mi}{1}\PYG{p}{]}\PYG{p}{]}\PYG{p}{]}
\end{sphinxVerbatim}

\end{fulllineitems}


\sphinxAtStartPar
Representation types of \(x\) correspond to a symplectic leaf of the quiver variety \(M_0(Q,x)\), and the symplectic leaf dimension is \(2p\) applied to the representation type, where \(p\) is the \sphinxcode{\sphinxupquote{p\_function}}.
\index{symplectic\_leaf\_dimension() (quivercombinatorics.quivercombinatorics.Quiver method)@\spxentry{symplectic\_leaf\_dimension()}\spxextra{quivercombinatorics.quivercombinatorics.Quiver method}}

\begin{fulllineitems}
\phantomsection\label{\detokenize{index:quivercombinatorics.quivercombinatorics.Quiver.symplectic_leaf_dimension}}
\pysigstartsignatures
\pysiglinewithargsret
{\sphinxcode{\sphinxupquote{Quiver.}}\sphinxbfcode{\sphinxupquote{symplectic\_leaf\_dimension}}}
{\sphinxparam{\DUrole{n}{tau}}}
{}
\pysigstopsignatures
\sphinxAtStartPar
Returns the dimension of the symplectic leaf corresponding to the representation type \(\tau\)

\sphinxAtStartPar
INPUT:
\begin{itemize}
\item {} 
\sphinxAtStartPar
\sphinxcode{\sphinxupquote{tau}} – a representation type, i.e., a list whose elements are 2\sphinxhyphen{}tuples, the first element is in \(\mathbb{N}Q_0\), and the second is in \(\mathbb{N}\)

\end{itemize}

\sphinxAtStartPar
OUTPUT: The corresponding symplectic leaf dimension, i.e., \(2p\) applied to every element of \(\tau\)

\sphinxAtStartPar
EXAMPLE:

\begin{sphinxVerbatim}[commandchars=\\\{\}]
\PYG{n}{sage}\PYG{p}{:} \PYG{k+kn}{from}\PYG{+w}{ }\PYG{n+nn}{quivercombinatorics}\PYG{+w}{ }\PYG{k+kn}{import} \PYG{o}{*}
\PYG{n}{sage}\PYG{p}{:} \PYG{n}{Q} \PYG{o}{=} \PYG{n}{Quiver}\PYG{p}{(}\PYG{p}{[}\PYG{p}{[}\PYG{l+m+mi}{0}\PYG{p}{,} \PYG{l+m+mi}{1}\PYG{p}{]}\PYG{p}{,} \PYG{p}{[}\PYG{l+m+mi}{1}\PYG{p}{,} \PYG{l+m+mi}{0}\PYG{p}{]}\PYG{p}{]}\PYG{p}{)}
\PYG{n}{sage}\PYG{p}{:} \PYG{n}{Q}\PYG{o}{.}\PYG{n}{symplectic\PYGZus{}leaf\PYGZus{}dimension}\PYG{p}{(}\PYG{p}{[}\PYG{p}{[}\PYG{p}{(}\PYG{l+m+mi}{1}\PYG{p}{,} \PYG{l+m+mi}{1}\PYG{p}{)}\PYG{p}{,} \PYG{l+m+mi}{2}\PYG{p}{]}\PYG{p}{,} \PYG{p}{[}\PYG{p}{(}\PYG{l+m+mi}{1}\PYG{p}{,} \PYG{l+m+mi}{1}\PYG{p}{)}\PYG{p}{,} \PYG{l+m+mi}{3}\PYG{p}{]}\PYG{p}{]}\PYG{p}{)}
\PYG{l+m+mi}{4}
\end{sphinxVerbatim}

\end{fulllineitems}

\index{CB\_decomposition() (quivercombinatorics.quivercombinatorics.Quiver method)@\spxentry{CB\_decomposition()}\spxextra{quivercombinatorics.quivercombinatorics.Quiver method}}

\begin{fulllineitems}
\phantomsection\label{\detokenize{index:quivercombinatorics.quivercombinatorics.Quiver.CB_decomposition}}
\pysigstartsignatures
\pysiglinewithargsret
{\sphinxcode{\sphinxupquote{Quiver.}}\sphinxbfcode{\sphinxupquote{CB\_decomposition}}}
{\sphinxparam{\DUrole{n}{l}}\sphinxparamcomma \sphinxparam{\DUrole{n}{x}}}
{}
\pysigstopsignatures
\sphinxAtStartPar
Returns the CB decomposition, which is the representation type maximizing \(p\)

\sphinxAtStartPar
INPUT:
\begin{itemize}
\item {} 
\sphinxAtStartPar
\sphinxcode{\sphinxupquote{l}} – an element of \(\mathbb{Z}Q_0\)

\item {} 
\sphinxAtStartPar
\sphinxcode{\sphinxupquote{x}} – an element of \(\mathbb{N}Q_0\)

\end{itemize}

\sphinxAtStartPar
OUTPUT: The CB decomposition with respect to \(x\), where \sphinxcode{\sphinxupquote{l}} is \(\lambda\)

\sphinxAtStartPar
EXAMPLE:

\begin{sphinxVerbatim}[commandchars=\\\{\}]
\PYG{n}{sage}\PYG{p}{:} \PYG{k+kn}{from}\PYG{+w}{ }\PYG{n+nn}{quivercombinatorics}\PYG{+w}{ }\PYG{k+kn}{import} \PYG{o}{*}
\PYG{n}{sage}\PYG{p}{:} \PYG{n}{Q} \PYG{o}{=} \PYG{n}{Quiver}\PYG{p}{(}\PYG{p}{[}\PYG{p}{[}\PYG{l+m+mi}{0}\PYG{p}{,} \PYG{l+m+mi}{1}\PYG{p}{]}\PYG{p}{,} \PYG{p}{[}\PYG{l+m+mi}{1}\PYG{p}{,} \PYG{l+m+mi}{0}\PYG{p}{]}\PYG{p}{]}\PYG{p}{)}
\PYG{n}{sage}\PYG{p}{:} \PYG{n}{Q}\PYG{o}{.}\PYG{n}{CB\PYGZus{}decomposition}\PYG{p}{(}\PYG{p}{(}\PYG{l+m+mi}{1}\PYG{p}{,} \PYG{o}{\PYGZhy{}}\PYG{l+m+mi}{1}\PYG{p}{)}\PYG{p}{,} \PYG{p}{(}\PYG{l+m+mi}{5}\PYG{p}{,} \PYG{l+m+mi}{5}\PYG{p}{)}\PYG{p}{)}
\PYG{p}{[}\PYG{p}{[}\PYG{p}{(}\PYG{l+m+mi}{1}\PYG{p}{,} \PYG{l+m+mi}{1}\PYG{p}{)}\PYG{p}{,} \PYG{l+m+mi}{1}\PYG{p}{]}\PYG{p}{,} \PYG{p}{[}\PYG{p}{(}\PYG{l+m+mi}{1}\PYG{p}{,} \PYG{l+m+mi}{1}\PYG{p}{)}\PYG{p}{,} \PYG{l+m+mi}{1}\PYG{p}{]}\PYG{p}{,} \PYG{p}{[}\PYG{p}{(}\PYG{l+m+mi}{1}\PYG{p}{,} \PYG{l+m+mi}{1}\PYG{p}{)}\PYG{p}{,} \PYG{l+m+mi}{1}\PYG{p}{]}\PYG{p}{,} \PYG{p}{[}\PYG{p}{(}\PYG{l+m+mi}{1}\PYG{p}{,} \PYG{l+m+mi}{1}\PYG{p}{)}\PYG{p}{,} \PYG{l+m+mi}{1}\PYG{p}{]}\PYG{p}{,} \PYG{p}{[}\PYG{p}{(}\PYG{l+m+mi}{1}\PYG{p}{,} \PYG{l+m+mi}{1}\PYG{p}{)}\PYG{p}{,} \PYG{l+m+mi}{1}\PYG{p}{]}\PYG{p}{]}
\end{sphinxVerbatim}

\end{fulllineitems}


\sphinxAtStartPar
Knowing the CB\sphinxhyphen{}decomposition, it is then easy to calculate the dimension of the quiver variety \(M_0(Q,x)\), as it gives the largest (open) symplectic leaf inside \(M_0(Q,x)\).
\index{quiver\_variety\_dimension() (quivercombinatorics.quivercombinatorics.Quiver method)@\spxentry{quiver\_variety\_dimension()}\spxextra{quivercombinatorics.quivercombinatorics.Quiver method}}

\begin{fulllineitems}
\phantomsection\label{\detokenize{index:quivercombinatorics.quivercombinatorics.Quiver.quiver_variety_dimension}}
\pysigstartsignatures
\pysiglinewithargsret
{\sphinxcode{\sphinxupquote{Quiver.}}\sphinxbfcode{\sphinxupquote{quiver\_variety\_dimension}}}
{\sphinxparam{\DUrole{n}{l}}\sphinxparamcomma \sphinxparam{\DUrole{n}{x}}}
{}
\pysigstopsignatures
\sphinxAtStartPar
Returns the dimension of the quiver variety

\sphinxAtStartPar
INPUT:
\begin{itemize}
\item {} 
\sphinxAtStartPar
\sphinxcode{\sphinxupquote{l}} – an element of \(\mathbb{Z}Q_0\)

\item {} 
\sphinxAtStartPar
\sphinxcode{\sphinxupquote{x}} – an element of \(\mathbb{N}Q_0\)

\end{itemize}

\sphinxAtStartPar
OUTPUT: The dimension of the corresponding quiver variety

\sphinxAtStartPar
EXAMPLE:

\begin{sphinxVerbatim}[commandchars=\\\{\}]
\PYG{n}{sage}\PYG{p}{:} \PYG{k+kn}{from}\PYG{+w}{ }\PYG{n+nn}{quivercombinatorics}\PYG{+w}{ }\PYG{k+kn}{import} \PYG{o}{*}
\PYG{n}{sage}\PYG{p}{:} \PYG{n}{Q} \PYG{o}{=} \PYG{n}{Quiver}\PYG{p}{(}\PYG{p}{[}\PYG{p}{[}\PYG{l+m+mi}{0}\PYG{p}{,} \PYG{l+m+mi}{1}\PYG{p}{]}\PYG{p}{,} \PYG{p}{[}\PYG{l+m+mi}{1}\PYG{p}{,} \PYG{l+m+mi}{0}\PYG{p}{]}\PYG{p}{]}\PYG{p}{)}
\PYG{n}{sage}\PYG{p}{:} \PYG{n}{Q}\PYG{o}{.}\PYG{n}{quiver\PYGZus{}variety\PYGZus{}dimension}\PYG{p}{(}\PYG{p}{(}\PYG{l+m+mi}{1}\PYG{p}{,} \PYG{o}{\PYGZhy{}}\PYG{l+m+mi}{1}\PYG{p}{)}\PYG{p}{,} \PYG{p}{(}\PYG{l+m+mi}{5}\PYG{p}{,} \PYG{l+m+mi}{5}\PYG{p}{)}\PYG{p}{)}
\PYG{l+m+mi}{10}
\end{sphinxVerbatim}

\end{fulllineitems}


\sphinxAtStartPar
It’s often important to know the codimension 2 leaves of the quiver variety \(M_0(Q,x)\), so we code this here.
\index{codimension\_two\_leaves() (quivercombinatorics.quivercombinatorics.Quiver method)@\spxentry{codimension\_two\_leaves()}\spxextra{quivercombinatorics.quivercombinatorics.Quiver method}}

\begin{fulllineitems}
\phantomsection\label{\detokenize{index:quivercombinatorics.quivercombinatorics.Quiver.codimension_two_leaves}}
\pysigstartsignatures
\pysiglinewithargsret
{\sphinxcode{\sphinxupquote{Quiver.}}\sphinxbfcode{\sphinxupquote{codimension\_two\_leaves}}}
{\sphinxparam{\DUrole{n}{l}}\sphinxparamcomma \sphinxparam{\DUrole{n}{x}}}
{}
\pysigstopsignatures
\sphinxAtStartPar
Returns the codimension 2 leaves of the quiver variety

\sphinxAtStartPar
INPUT:
\begin{itemize}
\item {} 
\sphinxAtStartPar
\sphinxcode{\sphinxupquote{l}} – an element of \(\mathbb{Z}Q_0\)

\item {} 
\sphinxAtStartPar
\sphinxcode{\sphinxupquote{x}} – an element of \(\mathbb{N}Q_0\)

\end{itemize}

\sphinxAtStartPar
OUTPUT: The representation types of the codimension 2 leaves of the quiver variety

\sphinxAtStartPar
EXAMPLE:

\begin{sphinxVerbatim}[commandchars=\\\{\}]
\PYG{n}{sage}\PYG{p}{:} \PYG{k+kn}{from}\PYG{+w}{ }\PYG{n+nn}{quivercombinatorics}\PYG{+w}{ }\PYG{k+kn}{import} \PYG{o}{*}
\PYG{n}{sage}\PYG{p}{:} \PYG{n}{Q} \PYG{o}{=} \PYG{n}{Quiver}\PYG{p}{(}\PYG{p}{[}\PYG{p}{[}\PYG{l+m+mi}{0}\PYG{p}{,} \PYG{l+m+mi}{1}\PYG{p}{]}\PYG{p}{,} \PYG{p}{[}\PYG{l+m+mi}{1}\PYG{p}{,} \PYG{l+m+mi}{0}\PYG{p}{]}\PYG{p}{]}\PYG{p}{)}
\PYG{n}{sage}\PYG{p}{:} \PYG{n}{Q}\PYG{o}{.}\PYG{n}{codimension\PYGZus{}two\PYGZus{}leaves}\PYG{p}{(}\PYG{p}{(}\PYG{l+m+mi}{0}\PYG{p}{,} \PYG{l+m+mi}{0}\PYG{p}{)}\PYG{p}{,} \PYG{p}{(}\PYG{l+m+mi}{5}\PYG{p}{,} \PYG{l+m+mi}{5}\PYG{p}{)}\PYG{p}{)}
\PYG{p}{[}\PYG{p}{[}\PYG{p}{[}\PYG{p}{(}\PYG{l+m+mi}{0}\PYG{p}{,} \PYG{l+m+mi}{1}\PYG{p}{)}\PYG{p}{,} \PYG{l+m+mi}{1}\PYG{p}{]}\PYG{p}{,}
\PYG{p}{[}\PYG{p}{(}\PYG{l+m+mi}{1}\PYG{p}{,} \PYG{l+m+mi}{0}\PYG{p}{)}\PYG{p}{,} \PYG{l+m+mi}{1}\PYG{p}{]}\PYG{p}{,}
\PYG{p}{[}\PYG{p}{(}\PYG{l+m+mi}{1}\PYG{p}{,} \PYG{l+m+mi}{1}\PYG{p}{)}\PYG{p}{,} \PYG{l+m+mi}{1}\PYG{p}{]}\PYG{p}{,}
\PYG{p}{[}\PYG{p}{(}\PYG{l+m+mi}{1}\PYG{p}{,} \PYG{l+m+mi}{1}\PYG{p}{)}\PYG{p}{,} \PYG{l+m+mi}{1}\PYG{p}{]}\PYG{p}{,}
\PYG{p}{[}\PYG{p}{(}\PYG{l+m+mi}{1}\PYG{p}{,} \PYG{l+m+mi}{1}\PYG{p}{)}\PYG{p}{,} \PYG{l+m+mi}{1}\PYG{p}{]}\PYG{p}{,}
\PYG{p}{[}\PYG{p}{(}\PYG{l+m+mi}{1}\PYG{p}{,} \PYG{l+m+mi}{1}\PYG{p}{)}\PYG{p}{,} \PYG{l+m+mi}{1}\PYG{p}{]}\PYG{p}{]}\PYG{p}{,}
\PYG{p}{[}\PYG{p}{[}\PYG{p}{(}\PYG{l+m+mi}{1}\PYG{p}{,} \PYG{l+m+mi}{1}\PYG{p}{)}\PYG{p}{,} \PYG{l+m+mi}{1}\PYG{p}{]}\PYG{p}{,} \PYG{p}{[}\PYG{p}{(}\PYG{l+m+mi}{1}\PYG{p}{,} \PYG{l+m+mi}{1}\PYG{p}{)}\PYG{p}{,} \PYG{l+m+mi}{1}\PYG{p}{]}\PYG{p}{,} \PYG{p}{[}\PYG{p}{(}\PYG{l+m+mi}{1}\PYG{p}{,} \PYG{l+m+mi}{1}\PYG{p}{)}\PYG{p}{,} \PYG{l+m+mi}{1}\PYG{p}{]}\PYG{p}{,} \PYG{p}{[}\PYG{p}{(}\PYG{l+m+mi}{1}\PYG{p}{,} \PYG{l+m+mi}{1}\PYG{p}{)}\PYG{p}{,} \PYG{l+m+mi}{2}\PYG{p}{]}\PYG{p}{]}\PYG{p}{]}
\end{sphinxVerbatim}

\end{fulllineitems}



\chapter{\protect\(\mathrm{ext}\protect\)\sphinxhyphen{}quivers}
\label{\detokenize{index:mathrm-ext-quivers}}
\sphinxAtStartPar
Let \(\tau=(\beta^{(1)},n_1;\dots;\beta^{(k)},n_k)\) be a representation type of \(Q\). The \(\mathrm{ext}\)\sphinxstyleemphasis{\sphinxhyphen{}quiver} \(\tilde{Q}\) associated to the symplectic leaf is the quiver with \(k\) vertices where the number of edges between vertex \(i\) and vertex \(j\) is \(-(\beta^{(i)},\beta^{(j)})\) and the number of loops at vertex \(i\) is \(p(\beta^{(i)})\). The corresponding dimension vector of \(\tilde{Q}\) is \(\mathbf{n}=(n_1,\dots,n_k)\), and we take \(\tilde{\lambda}=0\).
\index{ext\_quiver() (quivercombinatorics.quivercombinatorics.Quiver method)@\spxentry{ext\_quiver()}\spxextra{quivercombinatorics.quivercombinatorics.Quiver method}}

\begin{fulllineitems}
\phantomsection\label{\detokenize{index:quivercombinatorics.quivercombinatorics.Quiver.ext_quiver}}
\pysigstartsignatures
\pysiglinewithargsret
{\sphinxcode{\sphinxupquote{Quiver.}}\sphinxbfcode{\sphinxupquote{ext\_quiver}}}
{\sphinxparam{\DUrole{n}{tau}}}
{}
\pysigstopsignatures
\sphinxAtStartPar
Given a representation type, returns the ext\sphinxhyphen{}quiver

\sphinxAtStartPar
INPUT:
\begin{itemize}
\item {} 
\sphinxAtStartPar
\sphinxcode{\sphinxupquote{tau}} – a representation type, i.e., a list whose elements are 2\sphinxhyphen{}tuples, the first element is in \(\mathbb{N}Q_0\), and the second is in \(\mathbb{N}\)

\end{itemize}

\sphinxAtStartPar
OUTPUT: The \(\mathrm{ext}\)\sphinxhyphen{}quiver

\sphinxAtStartPar
EXAMPLE:

\begin{sphinxVerbatim}[commandchars=\\\{\}]
\PYG{n}{sage}\PYG{p}{:} \PYG{k+kn}{from}\PYG{+w}{ }\PYG{n+nn}{quivercombinatorics}\PYG{+w}{ }\PYG{k+kn}{import} \PYG{o}{*}
\PYG{n}{sage}\PYG{p}{:} \PYG{n}{Q} \PYG{o}{=} \PYG{n}{Quiver}\PYG{p}{(}\PYG{p}{[}\PYG{p}{[}\PYG{l+m+mi}{0}\PYG{p}{,} \PYG{l+m+mi}{1}\PYG{p}{]}\PYG{p}{,} \PYG{p}{[}\PYG{l+m+mi}{1}\PYG{p}{,} \PYG{l+m+mi}{0}\PYG{p}{]}\PYG{p}{]}\PYG{p}{)}
\PYG{n}{sage}\PYG{p}{:} \PYG{n}{tau} \PYG{o}{=} \PYG{p}{[}\PYG{p}{[}\PYG{p}{(}\PYG{l+m+mi}{1}\PYG{p}{,} \PYG{l+m+mi}{1}\PYG{p}{)}\PYG{p}{,} \PYG{l+m+mi}{2}\PYG{p}{]}\PYG{p}{,} \PYG{p}{[}\PYG{p}{(}\PYG{l+m+mi}{1}\PYG{p}{,} \PYG{l+m+mi}{1}\PYG{p}{)}\PYG{p}{,} \PYG{l+m+mi}{1}\PYG{p}{]}\PYG{p}{,} \PYG{p}{[}\PYG{p}{(}\PYG{l+m+mi}{1}\PYG{p}{,} \PYG{l+m+mi}{1}\PYG{p}{)}\PYG{p}{,} \PYG{l+m+mi}{1}\PYG{p}{]}\PYG{p}{,} \PYG{p}{[}\PYG{p}{(}\PYG{l+m+mi}{1}\PYG{p}{,} \PYG{l+m+mi}{1}\PYG{p}{)}\PYG{p}{,} \PYG{l+m+mi}{1}\PYG{p}{]}\PYG{p}{]}
\PYG{n}{sage}\PYG{p}{:} \PYG{n}{Q}\PYG{o}{.}\PYG{n}{ext\PYGZus{}quiver}\PYG{p}{(}\PYG{n}{tau}\PYG{p}{)}
\PYG{n}{a} \PYG{n}{quiver} \PYG{k}{with} \PYG{l+m+mi}{4} \PYG{n}{vertices} \PYG{o+ow}{and} \PYG{l+m+mi}{4} \PYG{n}{arrows}
\end{sphinxVerbatim}

\end{fulllineitems}

\index{ext\_dimension\_vector() (in module quivercombinatorics)@\spxentry{ext\_dimension\_vector()}\spxextra{in module quivercombinatorics}}

\begin{fulllineitems}
\phantomsection\label{\detokenize{index:quivercombinatorics.ext_dimension_vector}}
\pysigstartsignatures
\pysiglinewithargsret
{\sphinxcode{\sphinxupquote{quivercombinatorics.}}\sphinxbfcode{\sphinxupquote{ext\_dimension\_vector}}}
{\sphinxparam{\DUrole{n}{tau}}}
{}
\pysigstopsignatures
\sphinxAtStartPar
For a representation type \(\tau\), returns the associated dimension vector

\sphinxAtStartPar
INPUT:
\begin{itemize}
\item {} 
\sphinxAtStartPar
\sphinxcode{\sphinxupquote{tau}} – a representation type, i.e., a list whose elements are 2\sphinxhyphen{}tuples, the first element is in \(\mathbb{N}Q_0\), and the second is in \(\mathbb{N}\)

\end{itemize}

\sphinxAtStartPar
OUTPUT: a vector in \(\mathbb{N}Q_0\)

\sphinxAtStartPar
EXAMPLE:

\begin{sphinxVerbatim}[commandchars=\\\{\}]
\PYG{n}{sage}\PYG{p}{:} \PYG{k+kn}{from}\PYG{+w}{ }\PYG{n+nn}{quivercombinatorics}\PYG{+w}{ }\PYG{k+kn}{import} \PYG{o}{*}
\PYG{n}{sage}\PYG{p}{:} \PYG{n}{tau} \PYG{o}{=} \PYG{p}{[}\PYG{p}{[}\PYG{p}{(}\PYG{l+m+mi}{1}\PYG{p}{,} \PYG{l+m+mi}{1}\PYG{p}{)}\PYG{p}{,} \PYG{l+m+mi}{2}\PYG{p}{]}\PYG{p}{,} \PYG{p}{[}\PYG{p}{(}\PYG{l+m+mi}{1}\PYG{p}{,} \PYG{l+m+mi}{1}\PYG{p}{)}\PYG{p}{,} \PYG{l+m+mi}{1}\PYG{p}{]}\PYG{p}{,} \PYG{p}{[}\PYG{p}{(}\PYG{l+m+mi}{1}\PYG{p}{,} \PYG{l+m+mi}{1}\PYG{p}{)}\PYG{p}{,} \PYG{l+m+mi}{1}\PYG{p}{]}\PYG{p}{,} \PYG{p}{[}\PYG{p}{(}\PYG{l+m+mi}{1}\PYG{p}{,} \PYG{l+m+mi}{1}\PYG{p}{)}\PYG{p}{,} \PYG{l+m+mi}{1}\PYG{p}{]}\PYG{p}{]}
\PYG{n}{sage}\PYG{p}{:} \PYG{n}{ext\PYGZus{}dimension\PYGZus{}vector}\PYG{p}{(}\PYG{n}{tau}\PYG{p}{)}
\PYG{p}{(}\PYG{l+m+mi}{2}\PYG{p}{,} \PYG{l+m+mi}{1}\PYG{p}{,} \PYG{l+m+mi}{1}\PYG{p}{,} \PYG{l+m+mi}{1}\PYG{p}{)}
\end{sphinxVerbatim}

\end{fulllineitems}



\chapter{Classification of minimal degenerations of symplectic leaves}
\label{\detokenize{index:classification-of-minimal-degenerations-of-symplectic-leaves}}
\sphinxAtStartPar
For a quiver \(Q\), a positive root \(\alpha\) is \sphinxstyleemphasis{minimal imaginary} if it is imaginary and given any positive root \(\beta\), \(\alpha>\beta\) implies \(\beta\) is real.
\index{is\_minimal\_imaginary\_root() (quivercombinatorics.quivercombinatorics.Quiver method)@\spxentry{is\_minimal\_imaginary\_root()}\spxextra{quivercombinatorics.quivercombinatorics.Quiver method}}

\begin{fulllineitems}
\phantomsection\label{\detokenize{index:quivercombinatorics.quivercombinatorics.Quiver.is_minimal_imaginary_root}}
\pysigstartsignatures
\pysiglinewithargsret
{\sphinxcode{\sphinxupquote{Quiver.}}\sphinxbfcode{\sphinxupquote{is\_minimal\_imaginary\_root}}}
{\sphinxparam{\DUrole{n}{x}}}
{}
\pysigstopsignatures
\sphinxAtStartPar
Tests whether a vector is a vector is a minimal imaginary root

\sphinxAtStartPar
INPUT:
\begin{itemize}
\item {} 
\sphinxAtStartPar
\sphinxcode{\sphinxupquote{x}} – a vector in \(\mathbb{Z}Q_0\)

\end{itemize}

\sphinxAtStartPar
OUTPUT: \sphinxcode{\sphinxupquote{True}} if \sphinxcode{\sphinxupquote{x}} is a minimal imaginary root, and \sphinxcode{\sphinxupquote{False}} otherwise

\sphinxAtStartPar
EXAMPLE:

\begin{sphinxVerbatim}[commandchars=\\\{\}]
\PYG{n}{sage}\PYG{p}{:} \PYG{k+kn}{from}\PYG{+w}{ }\PYG{n+nn}{quivercombinatorics}\PYG{+w}{ }\PYG{k+kn}{import} \PYG{o}{*}
\PYG{n}{sage}\PYG{p}{:} \PYG{n}{Q} \PYG{o}{=} \PYG{n}{KroneckerQuiver}\PYG{p}{(}\PYG{l+m+mi}{2}\PYG{p}{)}
\PYG{n}{sage}\PYG{p}{:} \PYG{n}{Q}\PYG{o}{.}\PYG{n}{is\PYGZus{}minimal\PYGZus{}imaginary\PYGZus{}root}\PYG{p}{(}\PYG{p}{(}\PYG{l+m+mi}{1}\PYG{p}{,} \PYG{l+m+mi}{1}\PYG{p}{)}\PYG{p}{)}
\PYG{k+kc}{True}
\end{sphinxVerbatim}

\end{fulllineitems}

\index{all\_minimal\_imaginary\_positive\_roots() (quivercombinatorics.quivercombinatorics.Quiver method)@\spxentry{all\_minimal\_imaginary\_positive\_roots()}\spxextra{quivercombinatorics.quivercombinatorics.Quiver method}}

\begin{fulllineitems}
\phantomsection\label{\detokenize{index:quivercombinatorics.quivercombinatorics.Quiver.all_minimal_imaginary_positive_roots}}
\pysigstartsignatures
\pysiglinewithargsret
{\sphinxcode{\sphinxupquote{Quiver.}}\sphinxbfcode{\sphinxupquote{all\_minimal\_imaginary\_positive\_roots}}}
{\sphinxparam{\DUrole{n}{v}}}
{}
\pysigstopsignatures
\sphinxAtStartPar
Returns all minimal imaginary positive roots up to a bound \(v\)

\sphinxAtStartPar
INPUT:
\begin{itemize}
\item {} 
\sphinxAtStartPar
\sphinxcode{\sphinxupquote{v}} – a vector in \(\mathbb{N}Q_0\)

\end{itemize}

\sphinxAtStartPar
OUTPUT: A list of all minimal imaginary positive roots up to a bound \(v\)

\sphinxAtStartPar
EXAMPLE:

\begin{sphinxVerbatim}[commandchars=\\\{\}]
\PYG{n}{sage}\PYG{p}{:} \PYG{k+kn}{from}\PYG{+w}{ }\PYG{n+nn}{quivercombinatorics}\PYG{+w}{ }\PYG{k+kn}{import} \PYG{o}{*}
\PYG{n}{sage}\PYG{p}{:} \PYG{n}{Q} \PYG{o}{=} \PYG{n}{KroneckerQuiver}\PYG{p}{(}\PYG{l+m+mi}{2}\PYG{p}{)}
\PYG{n}{sage}\PYG{p}{:} \PYG{n}{Q}\PYG{o}{.}\PYG{n}{all\PYGZus{}minimal\PYGZus{}imaginary\PYGZus{}positive\PYGZus{}roots}\PYG{p}{(}\PYG{p}{(}\PYG{l+m+mi}{3}\PYG{p}{,} \PYG{l+m+mi}{3}\PYG{p}{)}\PYG{p}{)}
\PYG{p}{[}\PYG{p}{(}\PYG{l+m+mi}{1}\PYG{p}{,} \PYG{l+m+mi}{1}\PYG{p}{)}\PYG{p}{]}
\end{sphinxVerbatim}

\end{fulllineitems}


\sphinxAtStartPar
Given a quiver \(Q\), a \sphinxstyleemphasis{subquiver} \(T = (T_0, T_1)\) is specified by a subset \(T_0\subset Q_0\) of the vertices. Then, for \(i, j \in T_0 \subset Q_0\) the number of edges between vertex \(i\) and \(j\) in \(T\) equals the number of edges between vertex \(i\) and \(j\) in \(Q\). In particular, if \(i\in T_0\) then the number of loops at \(i\) in \(T\) is the same as the number of loops at \(i\) in \(Q\). The support of a dimension vector \(v\) is the subset \(T_0 \subset Q_0\) of all vertices \(i\) such that \(v_i\neq 0\). We think of the support of \(v\) as a subquiver of \(Q\).

\sphinxAtStartPar
We say that the subquiver \(T\) is of \sphinxstyleemphasis{affine ADE type} if it is one of the graphs appearing in the classification of simply\sphinxhyphen{}laced affine Dynkin diagrams \(\tilde{A}_l\), \(\tilde{D}_l\) or \(\tilde{E}_6\), \(\tilde{E}_7\), \(\tilde{E}_8\). Then there is a unique imaginary root \(\delta\), with \(p(\delta)=1\), whose support is \(T\). The key result is the classification of minimal imaginary roots is the following (Bellamy\sphinxhyphen{}Schedler 2026).

\begin{sphinxadmonition}{note}{Theorem (Bellamy\sphinxhyphen{}Schedler, 2026)}

\sphinxAtStartPar
If \(M\) is a closed leaf in \(M_0(Q,x)\) and \(M\subset\bar{L}\) is a minimal degeneration, then \(\bar{L}\cong M\times S\), where \(S\) is isomorphic to one of the following isolated singularities:
\begin{enumerate}
\sphinxsetlistlabels{\arabic}{enumi}{enumii}{}{.}%
\item {} 
\sphinxAtStartPar
A Kleinian singularity \((\mathbb{C}^2/\Gamma, 0)\).

\item {} 
\sphinxAtStartPar
The type A minimal nilpotent orbit closure \((\mathcal{O}_{\min},0)\) in \(\mathfrak{sl}_r(\mathbb{C})\).

\item {} 
\sphinxAtStartPar
\((\mathbb{C}^{2g}/\mathbb{Z}_2,0)\).

\item {} 
\sphinxAtStartPar
\(\operatorname{Spec}(\mathbb{C}[x\mid\deg x\geq 2])\).

\end{enumerate}
\end{sphinxadmonition}

\sphinxAtStartPar
We define:
\begin{quote}
\begin{equation*}
\begin{split}\tilde{\tau}(\beta,n):=(\beta,n;e_1,\alpha_1-n\beta_1;\dots;e_r,\alpha_r-n\beta_r)\end{split}
\end{equation*}\end{quote}

\sphinxAtStartPar
In each of the above cases, the leaf \(L\) from this theorem corresponds to representation types:
\begin{quote}
\begin{enumerate}
\sphinxsetlistlabels{\arabic}{enumi}{enumii}{}{.}%
\item {} 
\sphinxAtStartPar
\(\tilde{\tau}(\delta,n)\) for \(\delta\) a minimal imaginary root on an affine ADE subquiver.

\item {} 
\sphinxAtStartPar
\(\tilde{\tau}((1,1),n)\) for a two\sphinxhyphen{}vertex subquiver with \(t\geq 3\) edges between the vertices and no loops at the vertices.

\item {} \begin{equation*}
\begin{split}\tau=(e_i,a;e_i,a;e_1,\alpha_1;\dots;e_{i-1},\alpha_{i-1};e_{i+1},\alpha_{i+1};\dots)\end{split}
\end{equation*}
\end{enumerate}

\sphinxAtStartPar
where \(i\) is a vertex with \(g\geq 1\) loops and \(2a=\alpha_i\).
\begin{enumerate}
\sphinxsetlistlabels{\arabic}{enumi}{enumii}{}{.}%
\setcounter{enumi}{3}
\item {} \begin{equation*}
\begin{split}\tau=(e_i,a;e_i,b;e_1,\alpha_1;\dots;e_{i-1},\alpha_{i-1};e_{i+1},\alpha_{i+1};\dots)\end{split}
\end{equation*}
\end{enumerate}

\sphinxAtStartPar
where \(i\) is a vertex with \(g\geq 1\) loops and \(0<a\neq b<\alpha_i\) with \(a+b=\alpha_i\).
\end{quote}

\sphinxAtStartPar
We assign the following labels to each subminimal representation type (i.e., representation types corresponding to a minimal degeneration):
\begin{enumerate}
\sphinxsetlistlabels{\arabic}{enumi}{enumii}{}{.}%
\item {} 
\sphinxAtStartPar
\(A_l,D_l\) or \(E_6,E_7,E_8\) for affine Dynkin diagram of type \(\tilde{A}_l,\tilde{D}_l\) or \(\tilde{E}_6,\tilde{E}_7,\tilde{E}_8\).

\item {} 
\sphinxAtStartPar
\(a_{r-1}\).

\item {} 
\sphinxAtStartPar
\(c_g\).

\item {} 
\sphinxAtStartPar
\(m_g\).

\end{enumerate}
\index{all\_subminimal\_representation\_types() (quivercombinatorics.quivercombinatorics.Quiver method)@\spxentry{all\_subminimal\_representation\_types()}\spxextra{quivercombinatorics.quivercombinatorics.Quiver method}}

\begin{fulllineitems}
\phantomsection\label{\detokenize{index:quivercombinatorics.quivercombinatorics.Quiver.all_subminimal_representation_types}}
\pysigstartsignatures
\pysiglinewithargsret
{\sphinxcode{\sphinxupquote{Quiver.}}\sphinxbfcode{\sphinxupquote{all\_subminimal\_representation\_types}}}
{\sphinxparam{\DUrole{n}{v}}}
{}
\pysigstopsignatures
\sphinxAtStartPar
Returns all subminimal representation types up to a bound \(v\), and classifies them by affine Dynkin type, or \(a_{r-1}\), \(c_{g}\), \(m_{g}\)

\sphinxAtStartPar
INPUT:
\begin{itemize}
\item {} 
\sphinxAtStartPar
\sphinxcode{\sphinxupquote{v}} – a vector in \(\mathbb{N}Q_0\)

\end{itemize}

\sphinxAtStartPar
OUTPUT: A list of 2\sphinxhyphen{}tuples, the first element is a subminimal representation type up to a bound \(v\), the second element is its corresponding affine Dynkin type, or \(a_{r-1}\), \(c_{g}\), \(m_{g}\)

\sphinxAtStartPar
EXAMPLES:

\begin{sphinxVerbatim}[commandchars=\\\{\}]
\PYG{n}{sage}\PYG{p}{:} \PYG{k+kn}{from}\PYG{+w}{ }\PYG{n+nn}{quivercombinatorics}\PYG{+w}{ }\PYG{k+kn}{import} \PYG{o}{*}
\PYG{n}{sage}\PYG{p}{:} \PYG{n}{Q} \PYG{o}{=} \PYG{n}{CyclicQuiver}\PYG{p}{(}\PYG{l+m+mi}{3}\PYG{p}{)}
\PYG{n}{sage}\PYG{p}{:} \PYG{n}{Q}\PYG{o}{.}\PYG{n}{all\PYGZus{}subminimal\PYGZus{}representation\PYGZus{}types}\PYG{p}{(}\PYG{p}{(}\PYG{l+m+mi}{3}\PYG{p}{,} \PYG{l+m+mi}{3}\PYG{p}{,} \PYG{l+m+mi}{3}\PYG{p}{)}\PYG{p}{)}
\PYG{p}{[}\PYG{p}{(}\PYG{p}{[}\PYG{p}{[}\PYG{p}{(}\PYG{l+m+mi}{0}\PYG{p}{,} \PYG{l+m+mi}{0}\PYG{p}{,} \PYG{l+m+mi}{1}\PYG{p}{)}\PYG{p}{,} \PYG{l+m+mi}{2}\PYG{p}{]}\PYG{p}{,} \PYG{p}{[}\PYG{p}{(}\PYG{l+m+mi}{0}\PYG{p}{,} \PYG{l+m+mi}{1}\PYG{p}{,} \PYG{l+m+mi}{0}\PYG{p}{)}\PYG{p}{,} \PYG{l+m+mi}{2}\PYG{p}{]}\PYG{p}{,} \PYG{p}{[}\PYG{p}{(}\PYG{l+m+mi}{1}\PYG{p}{,} \PYG{l+m+mi}{0}\PYG{p}{,} \PYG{l+m+mi}{0}\PYG{p}{)}\PYG{p}{,} \PYG{l+m+mi}{2}\PYG{p}{]}\PYG{p}{,} \PYG{p}{[}\PYG{p}{(}\PYG{l+m+mi}{1}\PYG{p}{,} \PYG{l+m+mi}{1}\PYG{p}{,} \PYG{l+m+mi}{1}\PYG{p}{)}\PYG{p}{,} \PYG{l+m+mi}{1}\PYG{p}{]}\PYG{p}{]}\PYG{p}{,} \PYG{l+s+s1}{\PYGZsq{}}\PYG{l+s+s1}{A\PYGZus{}}\PYG{l+s+si}{\PYGZob{}2\PYGZcb{}}\PYG{l+s+s1}{\PYGZsq{}}\PYG{p}{)}\PYG{p}{,}
\PYG{p}{(}\PYG{p}{[}\PYG{p}{[}\PYG{p}{(}\PYG{l+m+mi}{0}\PYG{p}{,} \PYG{l+m+mi}{0}\PYG{p}{,} \PYG{l+m+mi}{1}\PYG{p}{)}\PYG{p}{,} \PYG{l+m+mi}{1}\PYG{p}{]}\PYG{p}{,} \PYG{p}{[}\PYG{p}{(}\PYG{l+m+mi}{0}\PYG{p}{,} \PYG{l+m+mi}{1}\PYG{p}{,} \PYG{l+m+mi}{0}\PYG{p}{)}\PYG{p}{,} \PYG{l+m+mi}{1}\PYG{p}{]}\PYG{p}{,} \PYG{p}{[}\PYG{p}{(}\PYG{l+m+mi}{1}\PYG{p}{,} \PYG{l+m+mi}{0}\PYG{p}{,} \PYG{l+m+mi}{0}\PYG{p}{)}\PYG{p}{,} \PYG{l+m+mi}{1}\PYG{p}{]}\PYG{p}{,} \PYG{p}{[}\PYG{p}{(}\PYG{l+m+mi}{1}\PYG{p}{,} \PYG{l+m+mi}{1}\PYG{p}{,} \PYG{l+m+mi}{1}\PYG{p}{)}\PYG{p}{,} \PYG{l+m+mi}{2}\PYG{p}{]}\PYG{p}{]}\PYG{p}{,} \PYG{l+s+s1}{\PYGZsq{}}\PYG{l+s+s1}{A\PYGZus{}}\PYG{l+s+si}{\PYGZob{}2\PYGZcb{}}\PYG{l+s+s1}{\PYGZsq{}}\PYG{p}{)}\PYG{p}{,}
\PYG{p}{(}\PYG{p}{[}\PYG{p}{[}\PYG{p}{(}\PYG{l+m+mi}{1}\PYG{p}{,} \PYG{l+m+mi}{1}\PYG{p}{,} \PYG{l+m+mi}{1}\PYG{p}{)}\PYG{p}{,} \PYG{l+m+mi}{3}\PYG{p}{]}\PYG{p}{]}\PYG{p}{,} \PYG{l+s+s1}{\PYGZsq{}}\PYG{l+s+s1}{A\PYGZus{}}\PYG{l+s+si}{\PYGZob{}2\PYGZcb{}}\PYG{l+s+s1}{\PYGZsq{}}\PYG{p}{)}\PYG{p}{]}

\PYG{n}{sage}\PYG{p}{:} \PYG{k+kn}{from}\PYG{+w}{ }\PYG{n+nn}{quivercombinatorics}\PYG{+w}{ }\PYG{k+kn}{import} \PYG{o}{*}
\PYG{n}{sage}\PYG{p}{:} \PYG{n}{Q} \PYG{o}{=} \PYG{n}{LoopQuiver}\PYG{p}{(}\PYG{l+m+mi}{3}\PYG{p}{)}
\PYG{n}{sage}\PYG{p}{:} \PYG{n}{Q}\PYG{o}{.}\PYG{n}{all\PYGZus{}subminimal\PYGZus{}representation\PYGZus{}types}\PYG{p}{(}\PYG{p}{(}\PYG{l+m+mi}{5}\PYG{p}{)}\PYG{p}{)}
\PYG{p}{[}\PYG{p}{(}\PYG{p}{[}\PYG{p}{[}\PYG{p}{(}\PYG{l+m+mi}{1}\PYG{p}{)}\PYG{p}{,} \PYG{l+m+mi}{1}\PYG{p}{]}\PYG{p}{,} \PYG{p}{[}\PYG{p}{(}\PYG{l+m+mi}{1}\PYG{p}{)}\PYG{p}{,} \PYG{l+m+mi}{4}\PYG{p}{]}\PYG{p}{]}\PYG{p}{,} \PYG{l+s+s1}{\PYGZsq{}}\PYG{l+s+s1}{m\PYGZus{}}\PYG{l+s+si}{\PYGZob{}3\PYGZcb{}}\PYG{l+s+s1}{\PYGZsq{}}\PYG{p}{)}\PYG{p}{,} \PYG{p}{(}\PYG{p}{[}\PYG{p}{[}\PYG{p}{(}\PYG{l+m+mi}{1}\PYG{p}{)}\PYG{p}{,} \PYG{l+m+mi}{2}\PYG{p}{]}\PYG{p}{,} \PYG{p}{[}\PYG{p}{(}\PYG{l+m+mi}{1}\PYG{p}{)}\PYG{p}{,} \PYG{l+m+mi}{3}\PYG{p}{]}\PYG{p}{]}\PYG{p}{,} \PYG{l+s+s1}{\PYGZsq{}}\PYG{l+s+s1}{m\PYGZus{}}\PYG{l+s+si}{\PYGZob{}3\PYGZcb{}}\PYG{l+s+s1}{\PYGZsq{}}\PYG{p}{)}\PYG{p}{]}

\PYG{n}{sage}\PYG{p}{:} \PYG{k+kn}{from}\PYG{+w}{ }\PYG{n+nn}{quivercombinatorics}\PYG{+w}{ }\PYG{k+kn}{import} \PYG{o}{*}
\PYG{n}{sage}\PYG{p}{:} \PYG{n}{A} \PYG{o}{=} \PYG{p}{[}\PYG{p}{[}\PYG{l+m+mi}{0}\PYG{p}{,} \PYG{l+m+mi}{1}\PYG{p}{,} \PYG{l+m+mi}{0}\PYG{p}{,} \PYG{l+m+mi}{3}\PYG{p}{]}\PYG{p}{,} \PYG{p}{[}\PYG{l+m+mi}{1}\PYG{p}{,} \PYG{l+m+mi}{0}\PYG{p}{,} \PYG{l+m+mi}{0}\PYG{p}{,} \PYG{l+m+mi}{0}\PYG{p}{]}\PYG{p}{,} \PYG{p}{[}\PYG{l+m+mi}{0}\PYG{p}{,} \PYG{l+m+mi}{0}\PYG{p}{,} \PYG{l+m+mi}{2}\PYG{p}{,} \PYG{l+m+mi}{0}\PYG{p}{]}\PYG{p}{,} \PYG{p}{[}\PYG{l+m+mi}{0}\PYG{p}{,} \PYG{l+m+mi}{0}\PYG{p}{,} \PYG{l+m+mi}{0}\PYG{p}{,} \PYG{l+m+mi}{0}\PYG{p}{]}\PYG{p}{]}
\PYG{n}{sage}\PYG{p}{:} \PYG{n}{Q} \PYG{o}{=} \PYG{n}{Quiver}\PYG{p}{(}\PYG{n}{A}\PYG{p}{)}
\PYG{n}{sage}\PYG{p}{:} \PYG{n}{Q}\PYG{o}{.}\PYG{n}{all\PYGZus{}subminimal\PYGZus{}representation\PYGZus{}types}\PYG{p}{(}\PYG{p}{(}\PYG{l+m+mi}{1}\PYG{p}{,} \PYG{l+m+mi}{1}\PYG{p}{,} \PYG{l+m+mi}{4}\PYG{p}{,} \PYG{l+m+mi}{1}\PYG{p}{)}\PYG{p}{)}
\PYG{p}{[}\PYG{p}{(}\PYG{p}{[}\PYG{p}{[}\PYG{p}{(}\PYG{l+m+mi}{0}\PYG{p}{,} \PYG{l+m+mi}{0}\PYG{p}{,} \PYG{l+m+mi}{0}\PYG{p}{,} \PYG{l+m+mi}{1}\PYG{p}{)}\PYG{p}{,} \PYG{l+m+mi}{1}\PYG{p}{]}\PYG{p}{,} \PYG{p}{[}\PYG{p}{(}\PYG{l+m+mi}{0}\PYG{p}{,} \PYG{l+m+mi}{0}\PYG{p}{,} \PYG{l+m+mi}{1}\PYG{p}{,} \PYG{l+m+mi}{0}\PYG{p}{)}\PYG{p}{,} \PYG{l+m+mi}{4}\PYG{p}{]}\PYG{p}{,} \PYG{p}{[}\PYG{p}{(}\PYG{l+m+mi}{1}\PYG{p}{,} \PYG{l+m+mi}{1}\PYG{p}{,} \PYG{l+m+mi}{0}\PYG{p}{,} \PYG{l+m+mi}{0}\PYG{p}{)}\PYG{p}{,} \PYG{l+m+mi}{1}\PYG{p}{]}\PYG{p}{]}\PYG{p}{,} \PYG{l+s+s1}{\PYGZsq{}}\PYG{l+s+s1}{A\PYGZus{}}\PYG{l+s+si}{\PYGZob{}1\PYGZcb{}}\PYG{l+s+s1}{\PYGZsq{}}\PYG{p}{)}\PYG{p}{,}
\PYG{p}{(}\PYG{p}{[}\PYG{p}{[}\PYG{p}{(}\PYG{l+m+mi}{0}\PYG{p}{,} \PYG{l+m+mi}{0}\PYG{p}{,} \PYG{l+m+mi}{1}\PYG{p}{,} \PYG{l+m+mi}{0}\PYG{p}{)}\PYG{p}{,} \PYG{l+m+mi}{4}\PYG{p}{]}\PYG{p}{,} \PYG{p}{[}\PYG{p}{(}\PYG{l+m+mi}{0}\PYG{p}{,} \PYG{l+m+mi}{1}\PYG{p}{,} \PYG{l+m+mi}{0}\PYG{p}{,} \PYG{l+m+mi}{0}\PYG{p}{)}\PYG{p}{,} \PYG{l+m+mi}{1}\PYG{p}{]}\PYG{p}{,} \PYG{p}{[}\PYG{p}{(}\PYG{l+m+mi}{1}\PYG{p}{,} \PYG{l+m+mi}{0}\PYG{p}{,} \PYG{l+m+mi}{0}\PYG{p}{,} \PYG{l+m+mi}{1}\PYG{p}{)}\PYG{p}{,} \PYG{l+m+mi}{1}\PYG{p}{]}\PYG{p}{]}\PYG{p}{,} \PYG{l+s+s1}{\PYGZsq{}}\PYG{l+s+s1}{a\PYGZus{}}\PYG{l+s+si}{\PYGZob{}2\PYGZcb{}}\PYG{l+s+s1}{\PYGZsq{}}\PYG{p}{)}\PYG{p}{,}
\PYG{p}{(}\PYG{p}{[}\PYG{p}{[}\PYG{p}{(}\PYG{l+m+mi}{0}\PYG{p}{,} \PYG{l+m+mi}{0}\PYG{p}{,} \PYG{l+m+mi}{0}\PYG{p}{,} \PYG{l+m+mi}{1}\PYG{p}{)}\PYG{p}{,} \PYG{l+m+mi}{1}\PYG{p}{]}\PYG{p}{,}
\PYG{p}{[}\PYG{p}{(}\PYG{l+m+mi}{0}\PYG{p}{,} \PYG{l+m+mi}{0}\PYG{p}{,} \PYG{l+m+mi}{1}\PYG{p}{,} \PYG{l+m+mi}{0}\PYG{p}{)}\PYG{p}{,} \PYG{l+m+mi}{1}\PYG{p}{]}\PYG{p}{,}
\PYG{p}{[}\PYG{p}{(}\PYG{l+m+mi}{0}\PYG{p}{,} \PYG{l+m+mi}{0}\PYG{p}{,} \PYG{l+m+mi}{1}\PYG{p}{,} \PYG{l+m+mi}{0}\PYG{p}{)}\PYG{p}{,} \PYG{l+m+mi}{3}\PYG{p}{]}\PYG{p}{,}
\PYG{p}{[}\PYG{p}{(}\PYG{l+m+mi}{0}\PYG{p}{,} \PYG{l+m+mi}{1}\PYG{p}{,} \PYG{l+m+mi}{0}\PYG{p}{,} \PYG{l+m+mi}{0}\PYG{p}{)}\PYG{p}{,} \PYG{l+m+mi}{1}\PYG{p}{]}\PYG{p}{,}
\PYG{p}{[}\PYG{p}{(}\PYG{l+m+mi}{1}\PYG{p}{,} \PYG{l+m+mi}{0}\PYG{p}{,} \PYG{l+m+mi}{0}\PYG{p}{,} \PYG{l+m+mi}{0}\PYG{p}{)}\PYG{p}{,} \PYG{l+m+mi}{1}\PYG{p}{]}\PYG{p}{]}\PYG{p}{,}
\PYG{l+s+s1}{\PYGZsq{}}\PYG{l+s+s1}{m\PYGZus{}}\PYG{l+s+si}{\PYGZob{}2\PYGZcb{}}\PYG{l+s+s1}{\PYGZsq{}}\PYG{p}{)}\PYG{p}{,}
\PYG{p}{(}\PYG{p}{[}\PYG{p}{[}\PYG{p}{(}\PYG{l+m+mi}{0}\PYG{p}{,} \PYG{l+m+mi}{0}\PYG{p}{,} \PYG{l+m+mi}{0}\PYG{p}{,} \PYG{l+m+mi}{1}\PYG{p}{)}\PYG{p}{,} \PYG{l+m+mi}{1}\PYG{p}{]}\PYG{p}{,}
\PYG{p}{[}\PYG{p}{(}\PYG{l+m+mi}{0}\PYG{p}{,} \PYG{l+m+mi}{0}\PYG{p}{,} \PYG{l+m+mi}{1}\PYG{p}{,} \PYG{l+m+mi}{0}\PYG{p}{)}\PYG{p}{,} \PYG{l+m+mi}{2}\PYG{p}{]}\PYG{p}{,}
\PYG{p}{[}\PYG{p}{(}\PYG{l+m+mi}{0}\PYG{p}{,} \PYG{l+m+mi}{0}\PYG{p}{,} \PYG{l+m+mi}{1}\PYG{p}{,} \PYG{l+m+mi}{0}\PYG{p}{)}\PYG{p}{,} \PYG{l+m+mi}{2}\PYG{p}{]}\PYG{p}{,}
\PYG{p}{[}\PYG{p}{(}\PYG{l+m+mi}{0}\PYG{p}{,} \PYG{l+m+mi}{1}\PYG{p}{,} \PYG{l+m+mi}{0}\PYG{p}{,} \PYG{l+m+mi}{0}\PYG{p}{)}\PYG{p}{,} \PYG{l+m+mi}{1}\PYG{p}{]}\PYG{p}{,}
\PYG{p}{[}\PYG{p}{(}\PYG{l+m+mi}{1}\PYG{p}{,} \PYG{l+m+mi}{0}\PYG{p}{,} \PYG{l+m+mi}{0}\PYG{p}{,} \PYG{l+m+mi}{0}\PYG{p}{)}\PYG{p}{,} \PYG{l+m+mi}{1}\PYG{p}{]}\PYG{p}{]}\PYG{p}{,}
\PYG{l+s+s1}{\PYGZsq{}}\PYG{l+s+s1}{c\PYGZus{}}\PYG{l+s+si}{\PYGZob{}2\PYGZcb{}}\PYG{l+s+s1}{\PYGZsq{}}\PYG{p}{)}\PYG{p}{]}
\end{sphinxVerbatim}

\end{fulllineitems}



\chapter{Plotting minimal degenerations of symplectic leaves in Hasse diagrams}
\label{\detokenize{index:plotting-minimal-degenerations-of-symplectic-leaves-in-hasse-diagrams}}
\sphinxAtStartPar
To compute these Hasse diagrams, make sure you have \sphinxcode{\sphinxupquote{dot2tex}} and \sphinxcode{\sphinxupquote{graphviz}} installed.

\begin{sphinxVerbatim}[commandchars=\\\{\}]
\PYG{n}{sage} \PYG{o}{\PYGZhy{}}\PYG{o}{\PYGZhy{}}\PYG{n}{pip} \PYG{n}{install} \PYG{n}{dot2tex}
\PYG{n}{sage} \PYG{o}{\PYGZhy{}}\PYG{o}{\PYGZhy{}}\PYG{n}{pip} \PYG{n}{install} \PYG{n}{graphviz}
\end{sphinxVerbatim}

\sphinxAtStartPar
We visualize the poset of symplectic leaves (equivalently of representation types) as a Hasse diagram: The vertices in this diagram are the representation types and the minimal degenerations, represented by an edge connecting two representation types, correspond to the situation where \(L_1\leq L_2\) (equivalently, \(\tau_1\leq\tau_2\)) but there does not exist \(L_3\) such that \(L_1 < L_3 < L_2\). We wish to compute the Hasse diagram, for which there are two methods.

\begin{sphinxadmonition}{note}{Method 1}

\sphinxAtStartPar
Traverse the Hasse diagram from the bottom, starting with the lowest leaves. When we reach a leaf \(L_i\), we compute the \(\mathrm{ext}\)\sphinxhyphen{}quiver associated to \(L_i\). On this \(\mathrm{ext}\)\sphinxhyphen{}quiver, we compute subminimal representation types of \(\mathbf{n}\). For each such representation type \(\tilde{\tau}=(\gamma^{(1)},m_1;\dots;\gamma^{(r)},m_r)\), define:
\begin{equation*}
\begin{split}D(\tilde{\tau}):=\left(D(\gamma^{(1)}),m_1;\dots;D(\gamma^{(r)}),m_r\right)\end{split}
\end{equation*}
\sphinxAtStartPar
By (Bellamy\sphinxhyphen{}Schedler 2026), it is known that \(D(\tilde{\tau})\) is a representation type for \(v\).

\sphinxAtStartPar
In this way, we have a rule that assigns to each subminimal representation type of \(\tilde{Q}\) a representation type \(D(\tilde{\tau}(\alpha))\) of \(v\). It is known that \(\tau < D(\tilde{\tau}(\alpha))\) is a minimal degeneration and they all occur in this way. Proceeding in this way allows one to build the Hasse diagram from the bottom up.
\end{sphinxadmonition}

\sphinxAtStartPar
For Method 1, we need the following helper functions.
\index{D\_map() (in module quivercombinatorics)@\spxentry{D\_map()}\spxextra{in module quivercombinatorics}}

\begin{fulllineitems}
\phantomsection\label{\detokenize{index:quivercombinatorics.D_map}}
\pysigstartsignatures
\pysiglinewithargsret
{\sphinxcode{\sphinxupquote{quivercombinatorics.}}\sphinxbfcode{\sphinxupquote{D\_map}}}
{\sphinxparam{\DUrole{n}{m}}\sphinxparamcomma \sphinxparam{\DUrole{n}{tau}}}
{}
\pysigstopsignatures
\sphinxAtStartPar
Evaluates a map \(D:\mathbb{Z}^k\to\mathbb{Z}^n\) from dimension vectors of the \(\mathrm{ext}\)\sphinxhyphen{}quiver to dimension vectors of the original quiver by \(D(m_1,\dots,m_k):=\sum_{i=1}^k m_i\beta^{(i)}\)

\sphinxAtStartPar
INPUT:
\begin{itemize}
\item {} 
\sphinxAtStartPar
\sphinxcode{\sphinxupquote{m}} – a vector in \(\mathbb{Z}\operatorname{ext}(Q)_0\)

\item {} 
\sphinxAtStartPar
\sphinxcode{\sphinxupquote{tau}} – a representation type of the original quiver, i.e., a list whose elements are 2\sphinxhyphen{}tuples, the first element is in \(\mathbb{N}Q_0\), and the second is in \(\mathbb{N}\)

\end{itemize}

\sphinxAtStartPar
OUTPUT: a vector in \(\mathbb{Z}Q_0\)

\sphinxAtStartPar
EXAMPLE:

\begin{sphinxVerbatim}[commandchars=\\\{\}]
\PYG{n}{sage}\PYG{p}{:} \PYG{k+kn}{from}\PYG{+w}{ }\PYG{n+nn}{quivercombinatorics}\PYG{+w}{ }\PYG{k+kn}{import} \PYG{o}{*}
\PYG{n}{sage}\PYG{p}{:} \PYG{n}{tau} \PYG{o}{=} \PYG{p}{[}\PYG{p}{[}\PYG{p}{(}\PYG{l+m+mi}{1}\PYG{p}{,} \PYG{l+m+mi}{1}\PYG{p}{)}\PYG{p}{,} \PYG{l+m+mi}{2}\PYG{p}{]}\PYG{p}{,} \PYG{p}{[}\PYG{p}{(}\PYG{l+m+mi}{1}\PYG{p}{,} \PYG{l+m+mi}{1}\PYG{p}{)}\PYG{p}{,} \PYG{l+m+mi}{1}\PYG{p}{]}\PYG{p}{,} \PYG{p}{[}\PYG{p}{(}\PYG{l+m+mi}{1}\PYG{p}{,} \PYG{l+m+mi}{1}\PYG{p}{)}\PYG{p}{,} \PYG{l+m+mi}{1}\PYG{p}{]}\PYG{p}{,} \PYG{p}{[}\PYG{p}{(}\PYG{l+m+mi}{1}\PYG{p}{,} \PYG{l+m+mi}{1}\PYG{p}{)}\PYG{p}{,} \PYG{l+m+mi}{1}\PYG{p}{]}\PYG{p}{]}
\PYG{n}{sage}\PYG{p}{:} \PYG{n}{D\PYGZus{}map}\PYG{p}{(}\PYG{p}{(}\PYG{l+m+mi}{1}\PYG{p}{,} \PYG{l+m+mi}{2}\PYG{p}{,} \PYG{l+m+mi}{8}\PYG{p}{,} \PYG{o}{\PYGZhy{}}\PYG{l+m+mi}{3}\PYG{p}{)}\PYG{p}{,} \PYG{n}{tau}\PYG{p}{)}
\PYG{p}{(}\PYG{l+m+mi}{8}\PYG{p}{,} \PYG{l+m+mi}{8}\PYG{p}{)}
\end{sphinxVerbatim}

\end{fulllineitems}

\index{D\_lifting() (in module quivercombinatorics)@\spxentry{D\_lifting()}\spxextra{in module quivercombinatorics}}

\begin{fulllineitems}
\phantomsection\label{\detokenize{index:quivercombinatorics.D_lifting}}
\pysigstartsignatures
\pysiglinewithargsret
{\sphinxcode{\sphinxupquote{quivercombinatorics.}}\sphinxbfcode{\sphinxupquote{D\_lifting}}}
{\sphinxparam{\DUrole{n}{tau}}\sphinxparamcomma \sphinxparam{\DUrole{n}{L}}}
{}
\pysigstopsignatures
\sphinxAtStartPar
Applies \sphinxcode{\sphinxupquote{D\_map}} to a representation type of the \(\mathrm{ext}\)\sphinxhyphen{}quiver

\sphinxAtStartPar
INPUT:
\begin{itemize}
\item {} 
\sphinxAtStartPar
\sphinxcode{\sphinxupquote{L}} – a representation type of the \(\mathrm{ext}\)\sphinxhyphen{}quiver, i.e., a list whose elements are 2\sphinxhyphen{}tuples, the first element is in \(\mathbb{N}\mathrm{ext}(Q)_0\), and the second is in \(\mathbb{N}\)

\item {} 
\sphinxAtStartPar
\sphinxcode{\sphinxupquote{tau}} – a representation type, i.e., a list whose elements are 2\sphinxhyphen{}tuples, the first element is in \(\mathbb{N}Q_0\), and the second is in \(\mathbb{N}\)

\end{itemize}

\sphinxAtStartPar
OUTPUT: a representation type of the \(\mathrm{ext}\)\sphinxhyphen{}quiver

\sphinxAtStartPar
EXAMPLE:

\begin{sphinxVerbatim}[commandchars=\\\{\}]
\PYG{n}{sage}\PYG{p}{:} \PYG{k+kn}{from}\PYG{+w}{ }\PYG{n+nn}{quivercombinatorics}\PYG{+w}{ }\PYG{k+kn}{import} \PYG{o}{*}
\PYG{n}{sage}\PYG{p}{:} \PYG{n}{tau} \PYG{o}{=} \PYG{p}{[}\PYG{p}{[}\PYG{p}{(}\PYG{l+m+mi}{1}\PYG{p}{,} \PYG{l+m+mi}{1}\PYG{p}{)}\PYG{p}{,} \PYG{l+m+mi}{2}\PYG{p}{]}\PYG{p}{,} \PYG{p}{[}\PYG{p}{(}\PYG{l+m+mi}{1}\PYG{p}{,} \PYG{l+m+mi}{1}\PYG{p}{)}\PYG{p}{,} \PYG{l+m+mi}{1}\PYG{p}{]}\PYG{p}{,} \PYG{p}{[}\PYG{p}{(}\PYG{l+m+mi}{1}\PYG{p}{,} \PYG{l+m+mi}{1}\PYG{p}{)}\PYG{p}{,} \PYG{l+m+mi}{1}\PYG{p}{]}\PYG{p}{,} \PYG{p}{[}\PYG{p}{(}\PYG{l+m+mi}{1}\PYG{p}{,} \PYG{l+m+mi}{1}\PYG{p}{)}\PYG{p}{,} \PYG{l+m+mi}{1}\PYG{p}{]}\PYG{p}{]}
\PYG{n}{sage}\PYG{p}{:} \PYG{n}{D\PYGZus{}lifting}\PYG{p}{(}\PYG{p}{[}\PYG{p}{[}\PYG{p}{(}\PYG{l+m+mi}{1}\PYG{p}{,} \PYG{l+m+mi}{2}\PYG{p}{,} \PYG{l+m+mi}{8}\PYG{p}{,} \PYG{o}{\PYGZhy{}}\PYG{l+m+mi}{3}\PYG{p}{)}\PYG{p}{,} \PYG{l+m+mi}{3}\PYG{p}{]}\PYG{p}{]}\PYG{p}{,} \PYG{n}{tau}\PYG{p}{)}
\PYG{p}{[}\PYG{p}{[}\PYG{p}{(}\PYG{l+m+mi}{8}\PYG{p}{,} \PYG{l+m+mi}{8}\PYG{p}{)}\PYG{p}{,} \PYG{l+m+mi}{3}\PYG{p}{]}\PYG{p}{]}
\end{sphinxVerbatim}

\end{fulllineitems}

\index{minimal\_degenerations() (quivercombinatorics.quivercombinatorics.Quiver method)@\spxentry{minimal\_degenerations()}\spxextra{quivercombinatorics.quivercombinatorics.Quiver method}}

\begin{fulllineitems}
\phantomsection\label{\detokenize{index:quivercombinatorics.quivercombinatorics.Quiver.minimal_degenerations}}
\pysigstartsignatures
\pysiglinewithargsret
{\sphinxcode{\sphinxupquote{Quiver.}}\sphinxbfcode{\sphinxupquote{minimal\_degenerations}}}
{\sphinxparam{\DUrole{n}{L}}}
{}
\pysigstopsignatures
\sphinxAtStartPar
Returns all minimal degenerations of a a symplectic leaf, corresponding to the representation type \(L\)

\sphinxAtStartPar
INPUT:
\begin{itemize}
\item {} 
\sphinxAtStartPar
\sphinxcode{\sphinxupquote{L}} – a representation type, i.e., a list whose elements are 2\sphinxhyphen{}tuples, the first element is in \(\mathbb{N}Q_0\), and the second is in \(\mathbb{N}\)

\end{itemize}

\sphinxAtStartPar
OUTPUT: A list of representation types

\sphinxAtStartPar
EXAMPLES:

\begin{sphinxVerbatim}[commandchars=\\\{\}]
\PYG{n}{sage}\PYG{p}{:} \PYG{k+kn}{from}\PYG{+w}{ }\PYG{n+nn}{quivercombinatorics}\PYG{+w}{ }\PYG{k+kn}{import} \PYG{o}{*}
\PYG{n}{sage}\PYG{p}{:} \PYG{n}{A} \PYG{o}{=} \PYG{p}{[}\PYG{p}{[}\PYG{l+m+mi}{0}\PYG{p}{,} \PYG{l+m+mi}{1}\PYG{p}{,} \PYG{l+m+mi}{0}\PYG{p}{,} \PYG{l+m+mi}{3}\PYG{p}{]}\PYG{p}{,} \PYG{p}{[}\PYG{l+m+mi}{1}\PYG{p}{,} \PYG{l+m+mi}{0}\PYG{p}{,} \PYG{l+m+mi}{0}\PYG{p}{,} \PYG{l+m+mi}{0}\PYG{p}{]}\PYG{p}{,} \PYG{p}{[}\PYG{l+m+mi}{0}\PYG{p}{,} \PYG{l+m+mi}{0}\PYG{p}{,} \PYG{l+m+mi}{2}\PYG{p}{,} \PYG{l+m+mi}{0}\PYG{p}{]}\PYG{p}{,} \PYG{p}{[}\PYG{l+m+mi}{0}\PYG{p}{,} \PYG{l+m+mi}{0}\PYG{p}{,} \PYG{l+m+mi}{0}\PYG{p}{,} \PYG{l+m+mi}{0}\PYG{p}{]}\PYG{p}{]}
\PYG{n}{sage}\PYG{p}{:} \PYG{n}{Q} \PYG{o}{=} \PYG{n}{Quiver}\PYG{p}{(}\PYG{n}{A}\PYG{p}{)}
\PYG{n}{sage}\PYG{p}{:} \PYG{n}{Q}\PYG{o}{.}\PYG{n}{minimal\PYGZus{}degenerations}\PYG{p}{(}\PYG{p}{[}\PYG{p}{[}\PYG{p}{(}\PYG{l+m+mi}{0}\PYG{p}{,} \PYG{l+m+mi}{0}\PYG{p}{,} \PYG{l+m+mi}{0}\PYG{p}{,} \PYG{l+m+mi}{1}\PYG{p}{)}\PYG{p}{,} \PYG{l+m+mi}{1}\PYG{p}{]}\PYG{p}{,} \PYG{p}{[}\PYG{p}{(}\PYG{l+m+mi}{0}\PYG{p}{,} \PYG{l+m+mi}{0}\PYG{p}{,} \PYG{l+m+mi}{1}\PYG{p}{,} \PYG{l+m+mi}{0}\PYG{p}{)}\PYG{p}{,} \PYG{l+m+mi}{4}\PYG{p}{]}\PYG{p}{,} \PYG{p}{[}\PYG{p}{(}\PYG{l+m+mi}{1}\PYG{p}{,} \PYG{l+m+mi}{1}\PYG{p}{,} \PYG{l+m+mi}{0}\PYG{p}{,} \PYG{l+m+mi}{0}\PYG{p}{)}\PYG{p}{,} \PYG{l+m+mi}{1}\PYG{p}{]}\PYG{p}{]}\PYG{p}{)}
\PYG{p}{[}\PYG{p}{(}\PYG{p}{[}\PYG{p}{[}\PYG{p}{(}\PYG{l+m+mi}{0}\PYG{p}{,} \PYG{l+m+mi}{0}\PYG{p}{,} \PYG{l+m+mi}{1}\PYG{p}{,} \PYG{l+m+mi}{0}\PYG{p}{)}\PYG{p}{,} \PYG{l+m+mi}{4}\PYG{p}{]}\PYG{p}{,} \PYG{p}{[}\PYG{p}{(}\PYG{l+m+mi}{1}\PYG{p}{,} \PYG{l+m+mi}{1}\PYG{p}{,} \PYG{l+m+mi}{0}\PYG{p}{,} \PYG{l+m+mi}{1}\PYG{p}{)}\PYG{p}{,} \PYG{l+m+mi}{1}\PYG{p}{]}\PYG{p}{]}\PYG{p}{,} \PYG{l+s+s1}{\PYGZsq{}}\PYG{l+s+s1}{\PYGZdl{}a\PYGZus{}}\PYG{l+s+si}{\PYGZob{}2\PYGZcb{}}\PYG{l+s+s1}{(1)\PYGZdl{}}\PYG{l+s+s1}{\PYGZsq{}}\PYG{p}{)}\PYG{p}{,}
\PYG{p}{(}\PYG{p}{[}\PYG{p}{[}\PYG{p}{(}\PYG{l+m+mi}{0}\PYG{p}{,} \PYG{l+m+mi}{0}\PYG{p}{,} \PYG{l+m+mi}{0}\PYG{p}{,} \PYG{l+m+mi}{1}\PYG{p}{)}\PYG{p}{,} \PYG{l+m+mi}{1}\PYG{p}{]}\PYG{p}{,} \PYG{p}{[}\PYG{p}{(}\PYG{l+m+mi}{0}\PYG{p}{,} \PYG{l+m+mi}{0}\PYG{p}{,} \PYG{l+m+mi}{1}\PYG{p}{,} \PYG{l+m+mi}{0}\PYG{p}{)}\PYG{p}{,} \PYG{l+m+mi}{1}\PYG{p}{]}\PYG{p}{,} \PYG{p}{[}\PYG{p}{(}\PYG{l+m+mi}{0}\PYG{p}{,} \PYG{l+m+mi}{0}\PYG{p}{,} \PYG{l+m+mi}{1}\PYG{p}{,} \PYG{l+m+mi}{0}\PYG{p}{)}\PYG{p}{,} \PYG{l+m+mi}{3}\PYG{p}{]}\PYG{p}{,} \PYG{p}{[}\PYG{p}{(}\PYG{l+m+mi}{1}\PYG{p}{,} \PYG{l+m+mi}{1}\PYG{p}{,} \PYG{l+m+mi}{0}\PYG{p}{,} \PYG{l+m+mi}{0}\PYG{p}{)}\PYG{p}{,} \PYG{l+m+mi}{1}\PYG{p}{]}\PYG{p}{]}\PYG{p}{,}
\PYG{l+s+s1}{\PYGZsq{}}\PYG{l+s+s1}{\PYGZdl{}m\PYGZus{}}\PYG{l+s+si}{\PYGZob{}2\PYGZcb{}}\PYG{l+s+s1}{(1)\PYGZdl{}}\PYG{l+s+s1}{\PYGZsq{}}\PYG{p}{)}\PYG{p}{,}
\PYG{p}{(}\PYG{p}{[}\PYG{p}{[}\PYG{p}{(}\PYG{l+m+mi}{0}\PYG{p}{,} \PYG{l+m+mi}{0}\PYG{p}{,} \PYG{l+m+mi}{0}\PYG{p}{,} \PYG{l+m+mi}{1}\PYG{p}{)}\PYG{p}{,} \PYG{l+m+mi}{1}\PYG{p}{]}\PYG{p}{,} \PYG{p}{[}\PYG{p}{(}\PYG{l+m+mi}{0}\PYG{p}{,} \PYG{l+m+mi}{0}\PYG{p}{,} \PYG{l+m+mi}{1}\PYG{p}{,} \PYG{l+m+mi}{0}\PYG{p}{)}\PYG{p}{,} \PYG{l+m+mi}{2}\PYG{p}{]}\PYG{p}{,} \PYG{p}{[}\PYG{p}{(}\PYG{l+m+mi}{0}\PYG{p}{,} \PYG{l+m+mi}{0}\PYG{p}{,} \PYG{l+m+mi}{1}\PYG{p}{,} \PYG{l+m+mi}{0}\PYG{p}{)}\PYG{p}{,} \PYG{l+m+mi}{2}\PYG{p}{]}\PYG{p}{,} \PYG{p}{[}\PYG{p}{(}\PYG{l+m+mi}{1}\PYG{p}{,} \PYG{l+m+mi}{1}\PYG{p}{,} \PYG{l+m+mi}{0}\PYG{p}{,} \PYG{l+m+mi}{0}\PYG{p}{)}\PYG{p}{,} \PYG{l+m+mi}{1}\PYG{p}{]}\PYG{p}{]}\PYG{p}{,}
\PYG{l+s+s1}{\PYGZsq{}}\PYG{l+s+s1}{\PYGZdl{}c\PYGZus{}}\PYG{l+s+si}{\PYGZob{}2\PYGZcb{}}\PYG{l+s+s1}{(1)\PYGZdl{}}\PYG{l+s+s1}{\PYGZsq{}}\PYG{p}{)}\PYG{p}{]}
\end{sphinxVerbatim}

\end{fulllineitems}


\sphinxAtStartPar
The second method to compute this Hasse diagram is as follows.

\begin{sphinxadmonition}{note}{Method 2}

\sphinxAtStartPar
Consider the larger set \(\mathcal{D}\) of all decompositions of a given dimension vector \(v\). Here, a decomposition is simply a way of writing \(v\) as a sum of smaller dimension vectors with multiplicities:
\begin{equation*}
\begin{split}\left(n_1,\beta^{(1)};\dots;n_k,\beta^{(k)}\right),\end{split}
\end{equation*}
\sphinxAtStartPar
where the \(\beta^{(i)}\) are dimension vectors satisfying \(\beta^{(i)}\leq v\). Following \sphinxhref{http://matrix.uantwerpen.be/lieven.lebruyn/b2hd-LeBruyn1988b.html}{this paper}, we say that:
\begin{equation*}
\begin{split}\varepsilon=(p_1,\gamma^{(1)};\dots;p_l,\gamma^{(l)})>\rho=(m_1,\alpha^{(1)};\dots;m_r,\alpha^{(r)})\end{split}
\end{equation*}
\sphinxAtStartPar
is a \sphinxstyleemphasis{direct successor} (\(\varepsilon\) is the successor of \(\rho\)) if either:
\begin{enumerate}
\sphinxsetlistlabels{\arabic}{enumi}{enumii}{}{.}%
\item {} 
\sphinxAtStartPar
\(l=r+1\) and \(\varepsilon\) can be ordered such that for all \(1\leq i\leq r-1\) we have \((m_i,\alpha^{(i)})=(p_i,\gamma^{(i)})\) and \(m_r=p_r=p_{r+1},\gamma^{(r)}+\gamma^{(r+1)}=\alpha^{(r)}\).

\item {} 
\sphinxAtStartPar
\(l=r-1\) and \(\varepsilon\) can be ordered such that for all \(1\leq i\leq r-2\) we have \((m_i,\alpha^{(i)})=(p_i,\gamma^{(i)})\) and \(\alpha^{(r)}=\alpha^{(r-1)}=\gamma^{(r-1)},m_{r-1}+m_r=p_{r-1}\).

\end{enumerate}
\end{sphinxadmonition}

\sphinxAtStartPar
The symplectic leaves form a subset of this set of decompositions, and by restriction, they inherit a sub\sphinxhyphen{}poset structure. It is shown in (Bellamy\sphinxhyphen{}Schedler, 2026) that this is the partial ordering given by leaf closure. We need the following helper functions for Method 2.
\index{all\_decompositions() (quivercombinatorics.quivercombinatorics.Quiver method)@\spxentry{all\_decompositions()}\spxextra{quivercombinatorics.quivercombinatorics.Quiver method}}

\begin{fulllineitems}
\phantomsection\label{\detokenize{index:quivercombinatorics.quivercombinatorics.Quiver.all_decompositions}}
\pysigstartsignatures
\pysiglinewithargsret
{\sphinxcode{\sphinxupquote{Quiver.}}\sphinxbfcode{\sphinxupquote{all\_decompositions}}}
{\sphinxparam{\DUrole{n}{v}}}
{}
\pysigstopsignatures
\sphinxAtStartPar
Constructs all decompositions of a given dimension vector \(v\)

\sphinxAtStartPar
INPUT:
\begin{itemize}
\item {} 
\sphinxAtStartPar
\sphinxcode{\sphinxupquote{v}} – an element of \(\mathbb{N}Q_0\)

\end{itemize}

\sphinxAtStartPar
OUTPUT: A list of decompositions of \sphinxtitleref{v}

\sphinxAtStartPar
EXAMPLE:

\begin{sphinxVerbatim}[commandchars=\\\{\}]
\PYG{n}{sage}\PYG{p}{:} \PYG{k+kn}{from}\PYG{+w}{ }\PYG{n+nn}{quivercombinatorics}\PYG{+w}{ }\PYG{k+kn}{import} \PYG{o}{*}
\PYG{n}{sage}\PYG{p}{:} \PYG{n}{A} \PYG{o}{=} \PYG{p}{[}\PYG{p}{[}\PYG{l+m+mi}{0}\PYG{p}{,} \PYG{l+m+mi}{1}\PYG{p}{,} \PYG{l+m+mi}{0}\PYG{p}{,} \PYG{l+m+mi}{3}\PYG{p}{]}\PYG{p}{,} \PYG{p}{[}\PYG{l+m+mi}{1}\PYG{p}{,} \PYG{l+m+mi}{0}\PYG{p}{,} \PYG{l+m+mi}{0}\PYG{p}{,} \PYG{l+m+mi}{0}\PYG{p}{]}\PYG{p}{,} \PYG{p}{[}\PYG{l+m+mi}{0}\PYG{p}{,} \PYG{l+m+mi}{0}\PYG{p}{,} \PYG{l+m+mi}{2}\PYG{p}{,} \PYG{l+m+mi}{0}\PYG{p}{]}\PYG{p}{,} \PYG{p}{[}\PYG{l+m+mi}{0}\PYG{p}{,} \PYG{l+m+mi}{0}\PYG{p}{,} \PYG{l+m+mi}{0}\PYG{p}{,} \PYG{l+m+mi}{0}\PYG{p}{]}\PYG{p}{]}
\PYG{n}{sage}\PYG{p}{:} \PYG{n}{Q} \PYG{o}{=} \PYG{n}{Quiver}\PYG{p}{(}\PYG{n}{A}\PYG{p}{)}
\PYG{n}{sage}\PYG{p}{:} \PYG{n}{Q}\PYG{o}{.}\PYG{n}{all\PYGZus{}decompositions}\PYG{p}{(}\PYG{p}{(}\PYG{l+m+mi}{1}\PYG{p}{,} \PYG{l+m+mi}{0}\PYG{p}{,} \PYG{l+m+mi}{2}\PYG{p}{,} \PYG{l+m+mi}{0}\PYG{p}{)}\PYG{p}{)}
\PYG{p}{[}\PYG{p}{[}\PYG{p}{[}\PYG{p}{(}\PYG{l+m+mi}{0}\PYG{p}{,} \PYG{l+m+mi}{0}\PYG{p}{,} \PYG{l+m+mi}{1}\PYG{p}{,} \PYG{l+m+mi}{0}\PYG{p}{)}\PYG{p}{,} \PYG{l+m+mi}{1}\PYG{p}{]}\PYG{p}{,} \PYG{p}{[}\PYG{p}{(}\PYG{l+m+mi}{0}\PYG{p}{,} \PYG{l+m+mi}{0}\PYG{p}{,} \PYG{l+m+mi}{1}\PYG{p}{,} \PYG{l+m+mi}{0}\PYG{p}{)}\PYG{p}{,} \PYG{l+m+mi}{1}\PYG{p}{]}\PYG{p}{,} \PYG{p}{[}\PYG{p}{(}\PYG{l+m+mi}{1}\PYG{p}{,} \PYG{l+m+mi}{0}\PYG{p}{,} \PYG{l+m+mi}{0}\PYG{p}{,} \PYG{l+m+mi}{0}\PYG{p}{)}\PYG{p}{,} \PYG{l+m+mi}{1}\PYG{p}{]}\PYG{p}{]}\PYG{p}{,}
\PYG{p}{[}\PYG{p}{[}\PYG{p}{(}\PYG{l+m+mi}{0}\PYG{p}{,} \PYG{l+m+mi}{0}\PYG{p}{,} \PYG{l+m+mi}{1}\PYG{p}{,} \PYG{l+m+mi}{0}\PYG{p}{)}\PYG{p}{,} \PYG{l+m+mi}{1}\PYG{p}{]}\PYG{p}{,} \PYG{p}{[}\PYG{p}{(}\PYG{l+m+mi}{1}\PYG{p}{,} \PYG{l+m+mi}{0}\PYG{p}{,} \PYG{l+m+mi}{1}\PYG{p}{,} \PYG{l+m+mi}{0}\PYG{p}{)}\PYG{p}{,} \PYG{l+m+mi}{1}\PYG{p}{]}\PYG{p}{]}\PYG{p}{,}
\PYG{p}{[}\PYG{p}{[}\PYG{p}{(}\PYG{l+m+mi}{0}\PYG{p}{,} \PYG{l+m+mi}{0}\PYG{p}{,} \PYG{l+m+mi}{1}\PYG{p}{,} \PYG{l+m+mi}{0}\PYG{p}{)}\PYG{p}{,} \PYG{l+m+mi}{2}\PYG{p}{]}\PYG{p}{,} \PYG{p}{[}\PYG{p}{(}\PYG{l+m+mi}{1}\PYG{p}{,} \PYG{l+m+mi}{0}\PYG{p}{,} \PYG{l+m+mi}{0}\PYG{p}{,} \PYG{l+m+mi}{0}\PYG{p}{)}\PYG{p}{,} \PYG{l+m+mi}{1}\PYG{p}{]}\PYG{p}{]}\PYG{p}{,}
\PYG{p}{[}\PYG{p}{[}\PYG{p}{(}\PYG{l+m+mi}{0}\PYG{p}{,} \PYG{l+m+mi}{0}\PYG{p}{,} \PYG{l+m+mi}{2}\PYG{p}{,} \PYG{l+m+mi}{0}\PYG{p}{)}\PYG{p}{,} \PYG{l+m+mi}{1}\PYG{p}{]}\PYG{p}{,} \PYG{p}{[}\PYG{p}{(}\PYG{l+m+mi}{1}\PYG{p}{,} \PYG{l+m+mi}{0}\PYG{p}{,} \PYG{l+m+mi}{0}\PYG{p}{,} \PYG{l+m+mi}{0}\PYG{p}{)}\PYG{p}{,} \PYG{l+m+mi}{1}\PYG{p}{]}\PYG{p}{]}\PYG{p}{,}
\PYG{p}{[}\PYG{p}{[}\PYG{p}{(}\PYG{l+m+mi}{1}\PYG{p}{,} \PYG{l+m+mi}{0}\PYG{p}{,} \PYG{l+m+mi}{2}\PYG{p}{,} \PYG{l+m+mi}{0}\PYG{p}{)}\PYG{p}{,} \PYG{l+m+mi}{1}\PYG{p}{]}\PYG{p}{]}\PYG{p}{]}
\end{sphinxVerbatim}

\end{fulllineitems}

\index{is\_direct\_successor() (in module quivercombinatorics)@\spxentry{is\_direct\_successor()}\spxextra{in module quivercombinatorics}}

\begin{fulllineitems}
\phantomsection\label{\detokenize{index:quivercombinatorics.is_direct_successor}}
\pysigstartsignatures
\pysiglinewithargsret
{\sphinxcode{\sphinxupquote{quivercombinatorics.}}\sphinxbfcode{\sphinxupquote{is\_direct\_successor}}}
{\sphinxparam{\DUrole{n}{epsilon}}\sphinxparamcomma \sphinxparam{\DUrole{n}{rho}}}
{}
\pysigstopsignatures
\sphinxAtStartPar
Verifies whether \sphinxcode{\sphinxupquote{epsilon}} is the direct successor of \sphinxcode{\sphinxupquote{rho}}

\sphinxAtStartPar
INPUT:
\begin{itemize}
\item {} 
\sphinxAtStartPar
\sphinxcode{\sphinxupquote{epsilon}} – an element of \(\mathbb{N}^n\)

\item {} 
\sphinxAtStartPar
\sphinxcode{\sphinxupquote{rho}} – an element of \(\mathbb{N}^n\)

\end{itemize}

\sphinxAtStartPar
OUTPUT: \sphinxcode{\sphinxupquote{True}} if \sphinxcode{\sphinxupquote{epsilon}} is the direct successor of \sphinxcode{\sphinxupquote{rho}}, and \sphinxcode{\sphinxupquote{False}} otherwise

\sphinxAtStartPar
EXAMPLE:

\begin{sphinxVerbatim}[commandchars=\\\{\}]
\PYG{n}{sage}\PYG{p}{:} \PYG{k+kn}{from}\PYG{+w}{ }\PYG{n+nn}{quivercombinatorics}\PYG{+w}{ }\PYG{k+kn}{import} \PYG{o}{*}
\PYG{n}{sage}\PYG{p}{:} \PYG{n}{rho} \PYG{o}{=} \PYG{p}{[}\PYG{p}{(}\PYG{p}{(}\PYG{l+m+mi}{1}\PYG{p}{,} \PYG{l+m+mi}{0}\PYG{p}{)}\PYG{p}{,} \PYG{l+m+mi}{2}\PYG{p}{)}\PYG{p}{,} \PYG{p}{(}\PYG{p}{(}\PYG{l+m+mi}{1}\PYG{p}{,} \PYG{l+m+mi}{0}\PYG{p}{)}\PYG{p}{,} \PYG{l+m+mi}{3}\PYG{p}{)}\PYG{p}{]}
\PYG{n}{sage}\PYG{p}{:} \PYG{n}{epsilon} \PYG{o}{=} \PYG{p}{[}\PYG{p}{(}\PYG{p}{(}\PYG{l+m+mi}{1}\PYG{p}{,} \PYG{l+m+mi}{0}\PYG{p}{)}\PYG{p}{,} \PYG{l+m+mi}{5}\PYG{p}{)}\PYG{p}{]}
\PYG{n}{sage}\PYG{p}{:} \PYG{n}{is\PYGZus{}direct\PYGZus{}successor}\PYG{p}{(}\PYG{n}{epsilon}\PYG{p}{,} \PYG{n}{rho}\PYG{p}{)}
\PYG{k+kc}{True}
\end{sphinxVerbatim}

\end{fulllineitems}


\sphinxAtStartPar
We are now able to obtain and plot the Hasse diagram for Methods 1 and 2. Method 1 seems to be much faster, so it is enabled by default.
\index{get\_Hasse\_diagram() (quivercombinatorics.quivercombinatorics.Quiver method)@\spxentry{get\_Hasse\_diagram()}\spxextra{quivercombinatorics.quivercombinatorics.Quiver method}}

\begin{fulllineitems}
\phantomsection\label{\detokenize{index:quivercombinatorics.quivercombinatorics.Quiver.get_Hasse_diagram}}
\pysigstartsignatures
\pysiglinewithargsret
{\sphinxcode{\sphinxupquote{Quiver.}}\sphinxbfcode{\sphinxupquote{get\_Hasse\_diagram}}}
{\sphinxparam{\DUrole{n}{l}}\sphinxparamcomma \sphinxparam{\DUrole{n}{v}}\sphinxparamcomma \sphinxparam{\DUrole{n}{method}\DUrole{o}{=}\DUrole{default_value}{1}}}
{}
\pysigstopsignatures
\sphinxAtStartPar
Applies Method 1 or Method 2 to obtain data for the Hasse diagram of minimal degenerations

\sphinxAtStartPar
INPUT:
\begin{itemize}
\item {} 
\sphinxAtStartPar
\sphinxcode{\sphinxupquote{l}} – an element of \(\mathbb{Z}Q_0\)

\item {} 
\sphinxAtStartPar
\sphinxcode{\sphinxupquote{v}} – an element of \(\mathbb{N}Q_0\)

\item {} 
\sphinxAtStartPar
\sphinxcode{\sphinxupquote{method}} – When set to \sphinxcode{\sphinxupquote{1}}, Method 1 will be used to obtain the Hasse diagram of minimal degenerations. When set to \sphinxcode{\sphinxupquote{2}}, Method 1 will be used to obtain the Hasse diagram of minimal degenerations.

\end{itemize}

\sphinxAtStartPar
OUTPUT:
\begin{itemize}
\item {} 
\sphinxAtStartPar
\sphinxcode{\sphinxupquote{leaves\_poset}} – the underlying poset of the Hasse diagram of minimal degenerations

\item {} 
\sphinxAtStartPar
\sphinxcode{\sphinxupquote{all\_leaves}} – a list of all symplectic leaves, with respect to their representation types, so a list of representation types, i.e., a list whose elements are 2\sphinxhyphen{}tuples, the first element is in \(\mathbb{N}Q_0\), and the second is in \(\mathbb{N}\)

\item {} 
\sphinxAtStartPar
\sphinxcode{\sphinxupquote{dimensions}} – a list of dimensions of all the symplectic leaves, using \sphinxcode{\sphinxupquote{symplectic\_leaf\_dimension}}

\item {} 
\sphinxAtStartPar
\sphinxcode{\sphinxupquote{edge\_labels}} – a list of 3\sphinxhyphen{}tuples, the first two entries of each tuple define the edge, and the last entry is the edge label of the corresponding minimal degeneration, classified by \sphinxcode{\sphinxupquote{all\_subminimal\_representation\_types}}

\end{itemize}

\sphinxAtStartPar
EXAMPLE:

\begin{sphinxVerbatim}[commandchars=\\\{\}]
\PYG{n}{sage}\PYG{p}{:} \PYG{k+kn}{from}\PYG{+w}{ }\PYG{n+nn}{quivercombinatorics}\PYG{+w}{ }\PYG{k+kn}{import} \PYG{o}{*}
\PYG{n}{sage}\PYG{p}{:} \PYG{n}{C} \PYG{o}{=} \PYG{p}{[}\PYG{p}{[}\PYG{l+m+mi}{2}\PYG{p}{,} \PYG{o}{\PYGZhy{}}\PYG{l+m+mi}{1}\PYG{p}{,} \PYG{l+m+mi}{0}\PYG{p}{]}\PYG{p}{,} \PYG{p}{[}\PYG{o}{\PYGZhy{}}\PYG{l+m+mi}{1}\PYG{p}{,} \PYG{l+m+mi}{2}\PYG{p}{,} \PYG{o}{\PYGZhy{}}\PYG{l+m+mi}{2}\PYG{p}{]}\PYG{p}{,} \PYG{p}{[}\PYG{l+m+mi}{0}\PYG{p}{,} \PYG{o}{\PYGZhy{}}\PYG{l+m+mi}{2}\PYG{p}{,} \PYG{l+m+mi}{2}\PYG{p}{]}\PYG{p}{]}
\PYG{n}{sage}\PYG{p}{:} \PYG{n}{Q} \PYG{o}{=} \PYG{n}{quiver\PYGZus{}from\PYGZus{}cartan\PYGZus{}matrix}\PYG{p}{(}\PYG{n}{C}\PYG{p}{)}
\PYG{n}{sage}\PYG{p}{:} \PYG{n}{Q}\PYG{o}{.}\PYG{n}{get\PYGZus{}Hasse\PYGZus{}diagram}\PYG{p}{(}\PYG{p}{(}\PYG{l+m+mi}{0}\PYG{p}{,} \PYG{l+m+mi}{0}\PYG{p}{,} \PYG{l+m+mi}{0}\PYG{p}{)}\PYG{p}{,} \PYG{p}{(}\PYG{l+m+mi}{2}\PYG{p}{,} \PYG{l+m+mi}{4}\PYG{p}{,} \PYG{l+m+mi}{3}\PYG{p}{)}\PYG{p}{)}
\PYG{p}{(}\PYG{n}{Finite} \PYG{n}{poset} \PYG{n}{containing} \PYG{l+m+mi}{8} \PYG{n}{elements}\PYG{p}{,}
\PYG{p}{[}\PYG{p}{[}\PYG{p}{[}\PYG{p}{(}\PYG{l+m+mi}{0}\PYG{p}{,} \PYG{l+m+mi}{0}\PYG{p}{,} \PYG{l+m+mi}{1}\PYG{p}{)}\PYG{p}{,} \PYG{l+m+mi}{1}\PYG{p}{]}\PYG{p}{,}
\PYG{p}{[}\PYG{p}{(}\PYG{l+m+mi}{0}\PYG{p}{,} \PYG{l+m+mi}{1}\PYG{p}{,} \PYG{l+m+mi}{0}\PYG{p}{)}\PYG{p}{,} \PYG{l+m+mi}{2}\PYG{p}{]}\PYG{p}{,}
\PYG{p}{[}\PYG{p}{(}\PYG{l+m+mi}{0}\PYG{p}{,} \PYG{l+m+mi}{1}\PYG{p}{,} \PYG{l+m+mi}{1}\PYG{p}{)}\PYG{p}{,} \PYG{l+m+mi}{1}\PYG{p}{]}\PYG{p}{,}
\PYG{p}{[}\PYG{p}{(}\PYG{l+m+mi}{0}\PYG{p}{,} \PYG{l+m+mi}{1}\PYG{p}{,} \PYG{l+m+mi}{1}\PYG{p}{)}\PYG{p}{,} \PYG{l+m+mi}{1}\PYG{p}{]}\PYG{p}{,}
\PYG{p}{[}\PYG{p}{(}\PYG{l+m+mi}{1}\PYG{p}{,} \PYG{l+m+mi}{0}\PYG{p}{,} \PYG{l+m+mi}{0}\PYG{p}{)}\PYG{p}{,} \PYG{l+m+mi}{2}\PYG{p}{]}\PYG{p}{]}\PYG{p}{,}
\PYG{p}{[}\PYG{p}{[}\PYG{p}{(}\PYG{l+m+mi}{0}\PYG{p}{,} \PYG{l+m+mi}{0}\PYG{p}{,} \PYG{l+m+mi}{1}\PYG{p}{)}\PYG{p}{,} \PYG{l+m+mi}{1}\PYG{p}{]}\PYG{p}{,} \PYG{p}{[}\PYG{p}{(}\PYG{l+m+mi}{0}\PYG{p}{,} \PYG{l+m+mi}{1}\PYG{p}{,} \PYG{l+m+mi}{0}\PYG{p}{)}\PYG{p}{,} \PYG{l+m+mi}{2}\PYG{p}{]}\PYG{p}{,} \PYG{p}{[}\PYG{p}{(}\PYG{l+m+mi}{0}\PYG{p}{,} \PYG{l+m+mi}{1}\PYG{p}{,} \PYG{l+m+mi}{1}\PYG{p}{)}\PYG{p}{,} \PYG{l+m+mi}{2}\PYG{p}{]}\PYG{p}{,} \PYG{p}{[}\PYG{p}{(}\PYG{l+m+mi}{1}\PYG{p}{,} \PYG{l+m+mi}{0}\PYG{p}{,} \PYG{l+m+mi}{0}\PYG{p}{)}\PYG{p}{,} \PYG{l+m+mi}{2}\PYG{p}{]}\PYG{p}{]}\PYG{p}{,}
\PYG{p}{[}\PYG{p}{[}\PYG{p}{(}\PYG{l+m+mi}{0}\PYG{p}{,} \PYG{l+m+mi}{0}\PYG{p}{,} \PYG{l+m+mi}{1}\PYG{p}{)}\PYG{p}{,} \PYG{l+m+mi}{2}\PYG{p}{]}\PYG{p}{,} \PYG{p}{[}\PYG{p}{(}\PYG{l+m+mi}{0}\PYG{p}{,} \PYG{l+m+mi}{1}\PYG{p}{,} \PYG{l+m+mi}{0}\PYG{p}{)}\PYG{p}{,} \PYG{l+m+mi}{3}\PYG{p}{]}\PYG{p}{,} \PYG{p}{[}\PYG{p}{(}\PYG{l+m+mi}{0}\PYG{p}{,} \PYG{l+m+mi}{1}\PYG{p}{,} \PYG{l+m+mi}{1}\PYG{p}{)}\PYG{p}{,} \PYG{l+m+mi}{1}\PYG{p}{]}\PYG{p}{,} \PYG{p}{[}\PYG{p}{(}\PYG{l+m+mi}{1}\PYG{p}{,} \PYG{l+m+mi}{0}\PYG{p}{,} \PYG{l+m+mi}{0}\PYG{p}{)}\PYG{p}{,} \PYG{l+m+mi}{2}\PYG{p}{]}\PYG{p}{]}\PYG{p}{,}
\PYG{p}{[}\PYG{p}{[}\PYG{p}{(}\PYG{l+m+mi}{0}\PYG{p}{,} \PYG{l+m+mi}{0}\PYG{p}{,} \PYG{l+m+mi}{1}\PYG{p}{)}\PYG{p}{,} \PYG{l+m+mi}{3}\PYG{p}{]}\PYG{p}{,} \PYG{p}{[}\PYG{p}{(}\PYG{l+m+mi}{0}\PYG{p}{,} \PYG{l+m+mi}{1}\PYG{p}{,} \PYG{l+m+mi}{0}\PYG{p}{)}\PYG{p}{,} \PYG{l+m+mi}{4}\PYG{p}{]}\PYG{p}{,} \PYG{p}{[}\PYG{p}{(}\PYG{l+m+mi}{1}\PYG{p}{,} \PYG{l+m+mi}{0}\PYG{p}{,} \PYG{l+m+mi}{0}\PYG{p}{)}\PYG{p}{,} \PYG{l+m+mi}{2}\PYG{p}{]}\PYG{p}{]}\PYG{p}{,}
\PYG{p}{[}\PYG{p}{[}\PYG{p}{(}\PYG{l+m+mi}{0}\PYG{p}{,} \PYG{l+m+mi}{1}\PYG{p}{,} \PYG{l+m+mi}{0}\PYG{p}{)}\PYG{p}{,} \PYG{l+m+mi}{1}\PYG{p}{]}\PYG{p}{,}
\PYG{p}{[}\PYG{p}{(}\PYG{l+m+mi}{0}\PYG{p}{,} \PYG{l+m+mi}{1}\PYG{p}{,} \PYG{l+m+mi}{1}\PYG{p}{)}\PYG{p}{,} \PYG{l+m+mi}{1}\PYG{p}{]}\PYG{p}{,}
\PYG{p}{[}\PYG{p}{(}\PYG{l+m+mi}{0}\PYG{p}{,} \PYG{l+m+mi}{1}\PYG{p}{,} \PYG{l+m+mi}{1}\PYG{p}{)}\PYG{p}{,} \PYG{l+m+mi}{1}\PYG{p}{]}\PYG{p}{,}
\PYG{p}{[}\PYG{p}{(}\PYG{l+m+mi}{0}\PYG{p}{,} \PYG{l+m+mi}{1}\PYG{p}{,} \PYG{l+m+mi}{1}\PYG{p}{)}\PYG{p}{,} \PYG{l+m+mi}{1}\PYG{p}{]}\PYG{p}{,}
\PYG{p}{[}\PYG{p}{(}\PYG{l+m+mi}{1}\PYG{p}{,} \PYG{l+m+mi}{0}\PYG{p}{,} \PYG{l+m+mi}{0}\PYG{p}{)}\PYG{p}{,} \PYG{l+m+mi}{2}\PYG{p}{]}\PYG{p}{]}\PYG{p}{,}
\PYG{p}{[}\PYG{p}{[}\PYG{p}{(}\PYG{l+m+mi}{0}\PYG{p}{,} \PYG{l+m+mi}{1}\PYG{p}{,} \PYG{l+m+mi}{0}\PYG{p}{)}\PYG{p}{,} \PYG{l+m+mi}{1}\PYG{p}{]}\PYG{p}{,} \PYG{p}{[}\PYG{p}{(}\PYG{l+m+mi}{0}\PYG{p}{,} \PYG{l+m+mi}{1}\PYG{p}{,} \PYG{l+m+mi}{1}\PYG{p}{)}\PYG{p}{,} \PYG{l+m+mi}{1}\PYG{p}{]}\PYG{p}{,} \PYG{p}{[}\PYG{p}{(}\PYG{l+m+mi}{0}\PYG{p}{,} \PYG{l+m+mi}{1}\PYG{p}{,} \PYG{l+m+mi}{1}\PYG{p}{)}\PYG{p}{,} \PYG{l+m+mi}{2}\PYG{p}{]}\PYG{p}{,} \PYG{p}{[}\PYG{p}{(}\PYG{l+m+mi}{1}\PYG{p}{,} \PYG{l+m+mi}{0}\PYG{p}{,} \PYG{l+m+mi}{0}\PYG{p}{)}\PYG{p}{,} \PYG{l+m+mi}{2}\PYG{p}{]}\PYG{p}{]}\PYG{p}{,}
\PYG{p}{[}\PYG{p}{[}\PYG{p}{(}\PYG{l+m+mi}{0}\PYG{p}{,} \PYG{l+m+mi}{1}\PYG{p}{,} \PYG{l+m+mi}{0}\PYG{p}{)}\PYG{p}{,} \PYG{l+m+mi}{1}\PYG{p}{]}\PYG{p}{,} \PYG{p}{[}\PYG{p}{(}\PYG{l+m+mi}{0}\PYG{p}{,} \PYG{l+m+mi}{1}\PYG{p}{,} \PYG{l+m+mi}{1}\PYG{p}{)}\PYG{p}{,} \PYG{l+m+mi}{3}\PYG{p}{]}\PYG{p}{,} \PYG{p}{[}\PYG{p}{(}\PYG{l+m+mi}{1}\PYG{p}{,} \PYG{l+m+mi}{0}\PYG{p}{,} \PYG{l+m+mi}{0}\PYG{p}{)}\PYG{p}{,} \PYG{l+m+mi}{2}\PYG{p}{]}\PYG{p}{]}\PYG{p}{,}
\PYG{p}{[}\PYG{p}{[}\PYG{p}{(}\PYG{l+m+mi}{2}\PYG{p}{,} \PYG{l+m+mi}{4}\PYG{p}{,} \PYG{l+m+mi}{3}\PYG{p}{)}\PYG{p}{,} \PYG{l+m+mi}{1}\PYG{p}{]}\PYG{p}{]}\PYG{p}{]}\PYG{p}{,}
\PYG{p}{[}\PYG{l+m+mi}{4}\PYG{p}{,} \PYG{l+m+mi}{2}\PYG{p}{,} \PYG{l+m+mi}{2}\PYG{p}{,} \PYG{l+m+mi}{0}\PYG{p}{,} \PYG{l+m+mi}{6}\PYG{p}{,} \PYG{l+m+mi}{4}\PYG{p}{,} \PYG{l+m+mi}{2}\PYG{p}{,} \PYG{l+m+mi}{8}\PYG{p}{]}\PYG{p}{,}
\PYG{p}{[}\PYG{p}{(}\PYG{l+m+mi}{0}\PYG{p}{,} \PYG{l+m+mi}{4}\PYG{p}{,} \PYG{l+s+s1}{\PYGZsq{}}\PYG{l+s+s1}{A\PYGZus{}}\PYG{l+s+si}{\PYGZob{}1\PYGZcb{}}\PYG{l+s+s1}{(1)}\PYG{l+s+s1}{\PYGZsq{}}\PYG{p}{)}\PYG{p}{,}
\PYG{p}{(}\PYG{l+m+mi}{1}\PYG{p}{,} \PYG{l+m+mi}{5}\PYG{p}{,} \PYG{l+s+s1}{\PYGZsq{}}\PYG{l+s+s1}{A\PYGZus{}}\PYG{l+s+si}{\PYGZob{}1\PYGZcb{}}\PYG{l+s+s1}{(1)}\PYG{l+s+s1}{\PYGZsq{}}\PYG{p}{)}\PYG{p}{,}
\PYG{p}{(}\PYG{l+m+mi}{1}\PYG{p}{,} \PYG{l+m+mi}{0}\PYG{p}{,} \PYG{l+s+s1}{\PYGZsq{}}\PYG{l+s+s1}{c\PYGZus{}}\PYG{l+s+si}{\PYGZob{}1\PYGZcb{}}\PYG{l+s+s1}{(1)}\PYG{l+s+s1}{\PYGZsq{}}\PYG{p}{)}\PYG{p}{,}
\PYG{p}{(}\PYG{l+m+mi}{2}\PYG{p}{,} \PYG{l+m+mi}{0}\PYG{p}{,} \PYG{l+s+s1}{\PYGZsq{}}\PYG{l+s+s1}{A\PYGZus{}}\PYG{l+s+si}{\PYGZob{}1\PYGZcb{}}\PYG{l+s+s1}{(1)}\PYG{l+s+s1}{\PYGZsq{}}\PYG{p}{)}\PYG{p}{,}
\PYG{p}{(}\PYG{l+m+mi}{2}\PYG{p}{,} \PYG{l+m+mi}{5}\PYG{p}{,} \PYG{l+s+s1}{\PYGZsq{}}\PYG{l+s+s1}{A\PYGZus{}}\PYG{l+s+si}{\PYGZob{}1\PYGZcb{}}\PYG{l+s+s1}{(1)}\PYG{l+s+s1}{\PYGZsq{}}\PYG{p}{)}\PYG{p}{,}
\PYG{p}{(}\PYG{l+m+mi}{3}\PYG{p}{,} \PYG{l+m+mi}{2}\PYG{p}{,} \PYG{l+s+s1}{\PYGZsq{}}\PYG{l+s+s1}{A\PYGZus{}}\PYG{l+s+si}{\PYGZob{}1\PYGZcb{}}\PYG{l+s+s1}{(1)}\PYG{l+s+s1}{\PYGZsq{}}\PYG{p}{)}\PYG{p}{,}
\PYG{p}{(}\PYG{l+m+mi}{3}\PYG{p}{,} \PYG{l+m+mi}{1}\PYG{p}{,} \PYG{l+s+s1}{\PYGZsq{}}\PYG{l+s+s1}{A\PYGZus{}}\PYG{l+s+si}{\PYGZob{}1\PYGZcb{}}\PYG{l+s+s1}{(1)}\PYG{l+s+s1}{\PYGZsq{}}\PYG{p}{)}\PYG{p}{,}
\PYG{p}{(}\PYG{l+m+mi}{3}\PYG{p}{,} \PYG{l+m+mi}{6}\PYG{p}{,} \PYG{l+s+s1}{\PYGZsq{}}\PYG{l+s+s1}{A\PYGZus{}}\PYG{l+s+si}{\PYGZob{}1\PYGZcb{}}\PYG{l+s+s1}{(1)}\PYG{l+s+s1}{\PYGZsq{}}\PYG{p}{)}\PYG{p}{,}
\PYG{p}{(}\PYG{l+m+mi}{4}\PYG{p}{,} \PYG{l+m+mi}{7}\PYG{p}{,} \PYG{l+s+s1}{\PYGZsq{}}\PYG{l+s+s1}{D\PYGZus{}}\PYG{l+s+si}{\PYGZob{}4\PYGZcb{}}\PYG{l+s+s1}{(1)}\PYG{l+s+s1}{\PYGZsq{}}\PYG{p}{)}\PYG{p}{,}
\PYG{p}{(}\PYG{l+m+mi}{5}\PYG{p}{,} \PYG{l+m+mi}{4}\PYG{p}{,} \PYG{l+s+s1}{\PYGZsq{}}\PYG{l+s+s1}{c\PYGZus{}}\PYG{l+s+si}{\PYGZob{}1\PYGZcb{}}\PYG{l+s+s1}{(1)}\PYG{l+s+s1}{\PYGZsq{}}\PYG{p}{)}\PYG{p}{,}
\PYG{p}{(}\PYG{l+m+mi}{6}\PYG{p}{,} \PYG{l+m+mi}{5}\PYG{p}{,} \PYG{l+s+s1}{\PYGZsq{}}\PYG{l+s+s1}{m\PYGZus{}}\PYG{l+s+si}{\PYGZob{}1\PYGZcb{}}\PYG{l+s+s1}{(1)}\PYG{l+s+s1}{\PYGZsq{}}\PYG{p}{)}\PYG{p}{]}\PYG{p}{)}
\end{sphinxVerbatim}

\end{fulllineitems}

\index{plot\_Hasse\_diagram() (quivercombinatorics.quivercombinatorics.Quiver method)@\spxentry{plot\_Hasse\_diagram()}\spxextra{quivercombinatorics.quivercombinatorics.Quiver method}}

\begin{fulllineitems}
\phantomsection\label{\detokenize{index:quivercombinatorics.quivercombinatorics.Quiver.plot_Hasse_diagram}}
\pysigstartsignatures
\pysiglinewithargsret
{\sphinxcode{\sphinxupquote{Quiver.}}\sphinxbfcode{\sphinxupquote{plot\_Hasse\_diagram}}}
{\sphinxparam{\DUrole{n}{l}}\sphinxparamcomma \sphinxparam{\DUrole{n}{v}}\sphinxparamcomma \sphinxparam{\DUrole{n}{method}\DUrole{o}{=}\DUrole{default_value}{1}}\sphinxparamcomma \sphinxparam{\DUrole{n}{format}\DUrole{o}{=}\DUrole{default_value}{'tikz'}}\sphinxparamcomma \sphinxparam{\DUrole{n}{filename}\DUrole{o}{=}\DUrole{default_value}{'output'}}}
{}
\pysigstopsignatures
\sphinxAtStartPar
Applies Method 1 or Method 2 to plot the Hasse diagram for minimal degenerations

\sphinxAtStartPar
INPUT:
\begin{itemize}
\item {} 
\sphinxAtStartPar
\sphinxcode{\sphinxupquote{l}} – an element of \(\mathbb{Z}Q_0\)

\item {} 
\sphinxAtStartPar
\sphinxcode{\sphinxupquote{v}} – an element of \(\mathbb{N}Q_0\)

\item {} 
\sphinxAtStartPar
\sphinxcode{\sphinxupquote{method}} – When set to \sphinxcode{\sphinxupquote{1}}, Method 1 will be used to obtain the Hasse diagram of minimal degenerations. When set to \sphinxcode{\sphinxupquote{2}}, Method 1 will be used to obtain the Hasse diagram of minimal degenerations.

\item {} 
\sphinxAtStartPar
\sphinxcode{\sphinxupquote{format}} – Takes values “tikz”, “dot”, and “sage”, depending on whether you want a \sphinxcode{\sphinxupquote{.tex}} file with \sphinxstyleemphasis{tikzpicture}, a \sphinxcode{\sphinxupquote{.dot}} file, or a sage Hasse diagram output

\item {} 
\sphinxAtStartPar
\sphinxcode{\sphinxupquote{filename}} – filename of output

\end{itemize}

\sphinxAtStartPar
OUTPUT:
\begin{itemize}
\item {} 
\sphinxAtStartPar
the Hasse diagram for minimal degenerations

\end{itemize}

\sphinxAtStartPar
EXAMPLE:

\begin{sphinxVerbatim}[commandchars=\\\{\}]
sage: C = [[2, \PYGZhy{}1, 0], [\PYGZhy{}1, 2, \PYGZhy{}2], [0, \PYGZhy{}2, 2]]
sage: Q = quiver\PYGZus{}from\PYGZus{}cartan\PYGZus{}matrix(C)
sage: Q.plot\PYGZus{}Hasse\PYGZus{}diagram((0, 0, 0), (2, 4, 3))
\PYGZbs{}documentclass[tikz]\PYGZob{}standalone\PYGZcb{}
\PYGZbs{}begin\PYGZob{}document\PYGZcb{}

\PYGZbs{}begin\PYGZob{}tikzpicture\PYGZcb{}[\PYGZgt{}=latex,line join=bevel,]
\PYGZpc{}\PYGZpc{}
\PYGZbs{}begin\PYGZob{}scope\PYGZcb{}
    \PYGZbs{}pgfsetstrokecolor\PYGZob{}black\PYGZcb{}
\PYGZhy{}\PYGZhy{}\PYGZhy{}
77 lines not printed (5057 characters in total).
Use print to see the full content.
\PYGZhy{}\PYGZhy{}\PYGZhy{}
    \PYGZbs{}draw (221.5bp,185.5bp) node \PYGZob{}\PYGZdl{}c\PYGZus{}\PYGZob{}1\PYGZcb{}(1)\PYGZdl{}\PYGZcb{};
    \PYGZbs{}draw [\PYGZhy{}\PYGZgt{}] (6) ..controls (293.59bp,110.99bp) and (289.95bp,119.79bp)  .. (284.0bp,126.0bp) .. controls (277.86bp,132.41bp) and (269.81bp,137.25bp)  .. (5);
    \PYGZbs{}draw (320.5bp,118.5bp) node \PYGZob{}\PYGZdl{}m\PYGZus{}\PYGZob{}1\PYGZcb{}(1)\PYGZdl{}\PYGZcb{};
\PYGZpc{}
\PYGZbs{}end\PYGZob{}tikzpicture\PYGZcb{}
\PYGZbs{}end\PYGZob{}document\PYGZcb{}
\end{sphinxVerbatim}

\end{fulllineitems}


\sphinxAtStartPar
The example above should export the following diagram:

\noindent{\hspace*{\fill}\sphinxincludegraphics[width=0.850\linewidth]{{output}.pdf}\hspace*{\fill}}

\renewcommand{\indexname}{Index}
\printindex
\end{document}